% !TeX root = ../../Book5.tex
% 纲要标题：## 附录 A　组合代数与反射群
% 纲要位置：工作记录/计划纲要.md，第 4910 行。
% * Coxeter 群及长度和 Bruhat 序
% * 对称函数、Young 表的权计数、插入取逆与格路径恒等式
% * 分拆双射、Gaussian 面积公式及有限域旗的规范坐标计数
% * Hecke 标准基的左右算子模型、Kazhdan–Lusztig 规范基递推、低秩胞腔及参数特化

\chapter{组合代数与反射群}
\label{b5:app:A}
\BookChapterNumber{b5:app:A}

\begin{chapterabstract}
本附录由 Coxeter 词与长度进入 Bruhat 序\footnote{布吕阿（François Georges René Bruhat，1929-2007），法国数学家，研究代数群及其表示；相关次序以他命名。}、Hecke 代数和 Kazhdan--Lusztig 规范基，同时用 Young 表、插入算法及有限计数连接组合方法与表示论。
\end{chapterabstract}
\begin{learninggoals}
\begin{itemize}
\item 在对称群中计算约化词、长度及低秩 Bruhat 序。
\item 用横条链、行插入与逆插入处理 Young 表计数。
\item 根据作用群和权重建立轨道计数式，使用分次生成函数及有限域计数。
\item 比较 Hecke 代数的两组基，核验低秩乘法和胞腔表示矩阵。
\end{itemize}
\end{learninggoals}

在 \(S_3\) 中，用相邻换位 \(s_1=(12)\)、\(s_2=(23)\) 书写同一个置换，可以得到 \(s_1s_2s_1\) 与 \(s_2s_1s_2\) 两个不同的最短字。若只把「词」当作群元素的名字，这种不唯一会令人困惑；若记录字长及允许的交换、编织关系，它反而揭示了置换之间的次序结构。一般反射群也有相似的词与长度问题。我们将从这些可手算的关系出发，再看 Young 表、Hecke 代数\footnote{赫克（Erich Hecke，1887-1947），德国数学家，在解析数论、模形式及相关算子理论方面作出重要贡献。}和规范基怎样把组合数据转为表示论的计算。

\ProseParagraphBreak
所需基础是有限群作用及本卷第6、7章的 Young 表、Schur 函数和 Pieri 规则。一般 Coxeter 词性质以根的符号变化和秩二反射计算给出证明概要；在此基础上，Hecke 标准基和 Kazhdan--Lusztig 规范基\footnote{卡日丹（\OriginalName{Дмитрий Александрович Каждан}，1946-），俄裔以色列数学家，研究表示论；卢斯蒂格（Gheorghe Lusztig，1946-），罗马尼亚裔美国数学家，研究代数群及几何表示论。两人共同引入这里的规范基与多项式。}的证明在本附录展开。所选计数结果均作证明，并给出秩一和 \(S_3\) 的具体代数计算。一般胞腔分类、规范基正性和 Schubert 簇的上同调理论可留待进一步学习。

% !TeX root = ../../Book5.tex
% 纲要标题：### A.1 Coxeter 群、长度与 Bruhat 序
% 纲要位置：工作记录/计划纲要.md，第 4917 行。
% 纲要细目（保留原定写作范围）：
%
% #### A.1.1 反射关系与 Coxeter 系统
% #### A.1.2 约化表达式与逆序数
% #### A.1.3 Bruhat 序与弱序的区别
%
\section{Coxeter 群、长度与 Bruhat 序}
\label{b5:sec:A01}

Weyl 群的元素可以写成简单反射组成的词，从而通过生成关系进行计算。Coxeter 群把这种计算抽象出来。这里先给出一般定义，再在对称群中证明长度公式，并用低秩例子区分两种常见的偏序。

\subsection{反射关系与 Coxeter 系统}
\label{b5:subsec:A01:01}

\begin{definition}\label{b5:def:coxeter-system}
	设 \(S\) 为有限集合，整数或符号 \(m_{st}\) 满足 \(m_{ss}=1\)、\(m_{st}=m_{ts}\in\{2,3,\ldots,\infty\}\) 对 \(s\ne t\) 成立。由生成元 \(S\) 及关系
	\[
		s^2=1,\qquad (st)^{m_{st}}=1\quad(m_{st}<\infty)
	\]
	定义的群 \(W\)，连同指定的生成集 \(S\)，称为\term[Coxeter-xitong]{Coxeter 系统}[Coxeter system]。\(m_{st}=\infty\) 表示不添加这一对的关系。与某个 \(s\in S\) 共轭的群元素称为该系统的反射。
\end{definition}
\symidx{\((W,S)\)：Coxeter 系统}

当 \(s,t\) 为对合时，关系 \((st)^m=1\) 等价于两边各有 \(m\) 项的辫关系
\[
	sts\cdots=tst\cdots.
\]
记两边的交替词为 \(A,B\)。无论 \(m\) 为奇数还是偶数，将 \(B\) 反向读取，均有 \(AB^{-1}=(st)^m\)。因此 \((st)^m=1\) 等价于 \(A=B\)；这里不要求 \(B\) 自身为对合。因此 \(m=2\) 表示交换，\(m=3\) 表示 \(sts=tst\)。秩二情形设 \(r=st\)，则 \(srs=r^{-1}\)，由 \(t=sr\) 和 \(sr=r^{-1}s\)，可把所有 \(s\) 移到右端并约去相邻的 \(s^2\)，故每个词都可化为 \(r^i\) 或 \(r^is\)；若 \(m<\infty\)，再令 \(0\le i<m\)，群至多有 \(2m\) 个元素。正 \(m\) 边形的旋转和反射实现这些不同元素，所以 \(m\ge3\) 时恰为阶 \(2m\) 的二面体群；\(m=2\) 时，两个生成元是相互独立的交换对合，给出 \(C_2\times C_2\)。若 \(m=\infty\)，在实直线上令 \(s(x)=-x\)、\(t(x)=1-x\)，则 \(r(x)=x-1\)。所有 \(r^i,r^is\)、\(i\in\mathbb Z\)，分别是不同的平移和反射，故均互异，群为无限二面体群。

\ProseParagraphBreak
对称群 \(S_n\) 的邻位换位 \(s_i=(i,i+1)\) 满足
\[
	s_i^2=1,\qquad s_is_j=s_js_i\ (|i-j|>1),\qquad s_is_{i+1}s_i=s_{i+1}s_is_{i+1}.
\]
它们给出 \(A_{n-1}\) 型 Coxeter 系统。为证明由这些关系定义的抽象群确实同构于 \(S_n\)，记该群为 \(W_n\)，其子群 \(H\) 由 \(s_1,\ldots,s_{n-2}\) 生成。置
\[
	r_n=1,\qquad r_j=s_{n-1}s_{n-2}\cdots s_j\quad(1\le j<n).
\]
集合 \(\bigcup_{j=1}^n Hr_j\) 对右乘每个 \(s_i\) 封闭：若 \(i<j-1\)，将 \(s_i\) 交换到左侧；若 \(i=j-1\) 或 \(i=j\)，结果分别是 \(r_{j-1}\) 或 \(r_{j+1}\)；若 \(i>j\)，先让末尾的 \(s_i\) 与 \(s_j,\ldots,s_{i-2}\) 交换，再将 \(s_i s_{i-1}s_i\) 换成 \(s_{i-1}s_i s_{i-1}\)。把新出现的 \(s_{i-1}\) 与左边的 \(s_{i+1},\ldots,s_{n-1}\) 交换，便得 \(r_js_i=s_{i-1}r_j\)。对 \(r_n=1\) 则直接处理 \(i<n-1\) 与 \(i=n-1\)。故它覆盖 \(W_n\)，于是归纳得到 \(|W_n|\le n!\)。另一方面，把生成元送到邻位换位得到满射 \(W_n\to S_n\)，而 \(|S_n|=n!\)，所以此映射是同构。

\subsection{约化表达式与逆序数}
\label{b5:subsec:A01:02}

\begin{definition}\label{b5:def:coxeter-length}
	设 \((W,S)\) 为 Coxeter 系统，\(w\in W\)。在所有满足 \(w=s_1\cdots s_r\)、\(s_i\in S\)、\(r\ge0\) 的表达式中，最小的 \(r\) 记为 \(\ell(w)\)。达到这个长度的表达式称为\term[yuehua-biaodashi]{约化表达式}[reduced expression]；空词表示单位元。
\end{definition}
\symidx{\(\ell(w)\)：Coxeter 系统的简单生成元长度}

这里的“长度”依赖指定的 \(S\)，不是任意群生成集给出的不变数量。因为生成元为对合，反转一个最短词给出 \(\ell(w^{-1})=\ell(w)\)。

\begin{proposition}\label{b5:prop:permutation-length-inversions}
	设 \(n\ge1\) 为整数，\(w\in S_n\)，用单行记号 \(\bigl(w(1),\ldots,w(n)\bigr)\) 表示，并用邻位换位 \(s_i=(i,i+1)\)、\(1\le i<n\)，定义长度 \(\ell\)。则
	\[
		\ell(w)=\operatorname{inv}(w) \coloneqq\#\bigl\{\,(i,j)\mid i<j,\ w(i)>w(j)\,\bigr\}.
	\]
	特别地，\(\ell(ws_i)=\ell(w)+1\) 当且仅当 \(w(i)<w(i+1)\)。
\end{proposition}
\begin{proof}
	右乘 \(s_i\) 交换相邻位置的值。除它们彼此形成的那一对外，它们与其他位置贡献的逆序总数不变，所以逆序数恰增减一。因此从单位元出发的任意表达式至少有 \(\operatorname{inv}(w)\) 项。

	若 \(w\ne1\)，它的单行记号不严格递增，故存在相邻下降 \(w(i)>w(i+1)\)。右乘 \(s_i\) 减少一个逆序。重复此过程，经过恰 \(\operatorname{inv}(w)\) 步得到单位元；反转这些换位就给出同长表达式。因此下界达到，并得到最后的增减判断。
\end{proof}

例如 \(w=3412\) 有四个逆序，
\[
	w=s_2s_1s_3s_2.
\]
该词有四项，因而约化。逆序数也立即给出最长置换 \(w_0=(n,n-1,\ldots,1)\) 的长度 \(n(n-1)/2\)。

\subsection{Bruhat 序与弱序的区别}
\label{b5:subsec:A01:03}

\begin{definition}\label{b5:def:bruhat-order}
	设 \((W,S)\) 为 Coxeter 系统，\(\ell\) 为关于 \(S\) 的长度函数，\(T=\{\,wsw^{-1}\mid w\in W,\ s\in S\,\}\)，\(u,v\in W\)。若存在链
	\[
		u=u_0,\ u_1,\ldots,u_r=v,\qquad u_{i+1}=u_it_i,\quad t_i\in T,\quad \ell(u_{i+1})>\ell(u_i),
	\]
	则记 \(u\le v\)，称为\term[Bruhat-xu]{Bruhat 序}[Bruhat order]；允许空链。若每步只许右乘简单生成元且长度增加一，则得到右弱序。
\end{definition}
\symidx{\(u\le v\)：Coxeter 系统中的 Bruhat 序}

反身性和传递性由链直接得到；非空链长度严格增加，所以同时 \(u\le v\) 与 \(v\le u\) 只能有 \(u=v\)，这证明它们都是偏序。弱序可描述为约化词的前缀关系：增加简单生成元的链拼成约化词；反过来，每个约化词的前缀仍约化，否则用更短词代换前缀会缩短整个词。

\ProseParagraphBreak
在 \(S_n\) 中，反射就是所有换位，不限于邻位换位。以 \(S_3=\langle s,t\rangle\) 为例，Bruhat 序按长度分层为
\[
	\{1\},\qquad\{s,t\},\qquad\{st,ts\},\qquad\{sts=tst\}.
\]
相邻两层之间所有可能的连接都存在。特别地 \(s\le ts\)，因为 \(s(sts)=ts\) 且长度从一增加到二；但在右弱序中没有 \(s\le ts\)，因为长度差为一，从 \(s\) 出发只能右乘一个简单生成元，得到 \(ss=1\) 或 \(st\)，都不是 \(ts\)。 具体说，Bruhat 序在中间两层的覆盖关系是
\[
	s\lessdot st,\qquad s\lessdot ts,\qquad t\lessdot st,\qquad t\lessdot ts.
\]
右弱序只保留其中 \(s\lessdot st\)、\(t\lessdot ts\)：它们分别由右乘 \(t\)、\(s\) 得到；另外两条需要右乘非简单反射 \(sts\)。两种序虽有相同的长度层次，覆盖边却不同。

\ProseParagraphBreak
Bruhat 序还可以从旗的退化中理解。设 \(k\) 为代数闭域，\(G=\operatorname{GL}_2(k)\)，\(B\) 为可逆上三角矩阵组成的子群。\(G/B\) 与 \(k^2\) 中的直线一一对应：矩阵 \(g\) 对应 \(g\langle e_1\rangle\)，而稳定 \(\langle e_1\rangle\) 的矩阵恰是 \(B\)。在射影直线中，\(B\) 固定 \(\langle e_1\rangle\)，其余直线唯一写为 \(\langle a e_1+e_2\rangle\)，\(a\in k\)，并构成一个轨道。这两个轨道分别对应 \(1,s\in S_2\)，记为 \(C_1,C_s\)。在 Zariski 拓扑中，\(C_s\) 是仿射直线，闭包是整条射影直线，因而
\[
	C_1\subseteq\overline{C_s},\qquad 1<s.
\]
这里的序关系记录了开轨道中哪些旗能够退化为边界上的旗。

\ProseParagraphBreak
这种退化也能看出 \(S_3\) 中右弱序所缺少的 \(s<ts\)。改在 \(G=\operatorname{GL}_3(k)\) 中取上三角群 \(B\)，令 \(s=(12)\)、\(t=(23)\)，按单行记号有 \(s=213\)、\(ts=312\)。一个完整旗由直线 \(L\) 及包含它的平面 \(P\) 给出。\(s\) 所对应轨道中的旗为
\[
	L=\langle e_2+a e_1\rangle, \qquad P=\langle e_1,e_2\rangle, \qquad a\in k;
\]
\(ts\) 所对应的轨道则由
\[
	L=\langle e_3+b e_1+c e_2\rangle, \qquad P=\langle e_1,e_3+c e_2\rangle, \qquad b,c\in k
\]
给出。上三角矩阵分别把标准置换旗送到这些旗，并且每个参数都由旗唯一确定。固定第一式中的 \(a\)，考虑
\[
	L_z=\langle e_2+a e_1+z e_3\rangle, \qquad P_z=\langle e_1,e_2+z e_3\rangle.
\]
对 \(z\ne0\)，将直线的向量除以 \(z\)，便得到第二式中 \(b=a/z\)、\(c=1/z\) 的旗；在 \(z=0\) 处则得到原来的 \(s\) 轨道中的旗。因此 \(C_s\subseteq\overline{C_{ts}}\)。这解释了为什么只记录右乘简单反射的弱序不足以描述这些轨道的闭包关系。\subsecref{b5:subsec:A03:03}将用相同的坐标计算有限域上的旗数。

\ProseParagraphBreak
后面构造 Hecke 代数需要三项一般词性质。
\begin{theorem}[交换性质与子词判据]\label{b5:thm:coxeter-word-properties}
	设 \((W,S)\) 为 Coxeter 系统，\(\ell\) 为关于 \(S\) 的长度函数，\(T\) 为该系统的反射集合，\(\le\) 为 Bruhat 序。则对 \(u,v,w\in W\) 有下列性质。
	\begin{enumerate}
		\item 若 \(w=s_1\cdots s_r\) 为约化词，\(s_i\in S\)，反射 \(t\in T\) 满足 \(\ell(wt)<\ell(w)\)，则 \(wt\) 可由该词删去一个生成元得到。特别地，\(\ell(ws)=\ell(w)\pm1\) 对 \(s\in S\) 成立：下降时，删除一位及 \(\ell(w)\le\ell(ws)+1\) 共同给出差为一；上升时，将这一论证用于 \(ws\) 右乘 \(s\) 回到 \(w\)。
		\item 同一元素的任意两个约化词可用有限次辫关系相互变换。
		\item 对 \(v\) 的任意固定约化词，\(u\le v\) 当且仅当 \(u\) 是该词的一个子词的乘积。每个区间 \(\{u:u\le v\}\) 都有限。
	\end{enumerate}
\end{theorem}
\begin{proofsketch}
	在实向量空间 \(\bigoplus_{s\in S}\mathbb R\alpha_s\) 上规定
	\[
		B(\alpha_s,\alpha_s)=1,\qquad B(\alpha_s,\alpha_t)=-\cos(\pi/m_{st}),\qquad s(x)=x-2B(x,\alpha_s)\alpha_s,
	\]
	其中 \(m_{st}=\infty\) 时取 \(\cos(\pi/m_{st})=1\)。秩二反射计算验证这些算子满足群的生成关系。根是各个 \(w\alpha_s\)。基本的根符号引理说：每个根的简单根系数全非负或全非正，而且 \(w\alpha_s\) 为负根当且仅当 \(\ell(ws)<\ell(w)\)。其证明按词长归纳，将末尾两个不同生成元组成的交替词单独处理；有限二面体情形是相邻反射轴的夹角 \(\pi/m\) 的计算，无限情形的系数递推一直留在正锥中。这是本概要省略的主要归纳，详见\cite[Sections 5.3--5.7, pp. 108--116]{Humphreys1990ReflectionGroups}。

	这个输入将词长下降变成根的符号变化，从而识别应当删去的字母。简单反射只把它对应的简单正根变负，其余正根仍为正。根符号判别对一般反射给出：若 \(\ell(wt)<\ell(w)\)，则相应正根 \(\alpha\) 满足 \(w\alpha<0\)。写约化词 \(w=s_1\cdots s_r\)，记
	\[
		\beta_{r+1}=\alpha,\qquad \beta_k=s_k\cdots s_r\alpha\quad(1\le k\le r).
	\]
	从右端逐次作用时，起点 \(\beta_{r+1}>0\)，终点 \(\beta_1<0\)，故有第一次变号的位置 \(j\)，满足 \(\beta_{j+1}>0\)、\(s_j\beta_{j+1}<0\)。简单反射的上述性质迫使 \(\beta_{j+1}=\alpha_{s_j}\)。令 \(a=s_1\cdots s_{j-1}\)、\(b=s_{j+1}\cdots s_r\)，就有 \(w=a s_jb\) 及 \(b\alpha=\alpha_{s_j}\)，所以对应根反射满足 \(t=b^{-1}s_jb\)。于是
	\[
		wt=a s_j b\,b^{-1}s_jb=ab,
	\]
	恰好删去了位置 \(j\) 的字母。这给出第一项。对非约化词取最早的非约化前缀，设它为 \(us\)，其中 \(u\) 为约化前缀。若 \(\ell(us)=\ell(u)+1\)，该前缀便约化，与取法矛盾；由长度只差一可知 \(\ell(us)=\ell(u)-1\)。因此末字母使前一约化前缀变短；交换性质删去前缀中的一位，再删去末字母，得到不改变乘积的双位置删除。重复便可将任何非约化词变为约化子词。

	第二项按 \(\ell(w)\) 归纳。难点是两个约化词可能以不同的 \(s,t\) 开始；需要先将这两个首字母的差异集中到秩二子群中。为此使用标准抛物子群与分解结论：对 \(I\subseteq S\)，记 \(W_I=\langle I\rangle\)，则 \((W_I,I)\) 是具有原参数 \(m_{st}\) 的 Coxeter 系统，而且 \(W_I\) 内关于 \(I\) 的长度等于全群关于 \(S\) 的长度，见\cite[Section 5.5, pp. 113--114]{Humphreys1990ReflectionGroups}。每个陪集 \(W_Iw\) 有唯一的最短代表 \(z\)，并且
	\[
		w=w_Iz,\qquad w_I\in W_I,\qquad \ell(az)=\ell(a)+\ell(z)\quad(a\in W_I).
	\]
	因而分解中的两个因子唯一，且 \(\ell(w)=\ell(w_I)+\ell(z)\)。这一结论的存在性、唯一性和长度可加性均采用\cite[Section 5.12, p. 123]{Humphreys1990ReflectionGroups}中的抛物分解；原书写成陪集 \(zW_I\) 的形式，取逆便得到这里的形式。

	两个约化词首字母相同，就约去首字母；不同，记为 \(s,t\)，并取 \(I=\{s,t\}\)。由于 \(s,t\) 都使左乘 \(w\) 的长度下降，长度可加性说明它们也都使左乘 \(w_I\) 的长度下降。秩二群中只有有限二面体群的最长元具有这两个下降，故 \(W_I\) 有限且 \(w_I\) 为其最长元。它的两个交替约化词正由一条辫关系联系；将 \(z\) 的约化词接在后面，便得到分别以 \(s,t\) 开始的两个约化词。各自与原词比较时，约去共同首字母后可用归纳假设，再用这条辫关系把两者联系起来。

	在 \(A_2\) 型的 \(S_3\) 中，这一步就是 \(w_0=sts=tst\)，此时 \(w_I=w_0\)、\(z=1\)。一般的抛物分解将同一个秩二机制从较长的词中分离出来；长度可加性保证接回 \(z\) 时不会发生消去。

	若 \(u\le v\)，沿反射下降链重复使用第一项，每步所得词若不约化，就先作双位置删除；整个过程始终只从原词删除位置，故得到子词。反向的困难是：选取子词的末字母，对应的乘积未必增长。因此按原词长度归纳时，不选末字母可直接使用较短前缀；选末字母则使用提升性质。它根据 \(xs\) 的长度增减，将已知的 \(x\le v\) 转成 \(xs\le v\) 或 \(xs\le vs\)，从而把按反射定义的序关系与添入子词末字母的操作联系起来。提升性质的两个长度情形及此项归纳采用\cite[Sections 5.8--5.10, pp. 117--120]{Humphreys1990ReflectionGroups}。最后，一个长度 \(r\) 的固定词仅有 \(2^r\) 个位置子集，所以区间有限。
\end{proofsketch}

% !TeX root = ../../Book5.tex
% 纲要标题：### A.2 对称函数、Young 表与组合恒等式
% 纲要位置：工作记录/计划纲要.md，第 4925 行。
% 纲要细目（保留原定写作范围）：
%
% #### A.2.1 半标准表与 Schur 多项式
% #### A.2.2 行插入与成对标准表
% #### A.2.3 维数恒等式与格路径计数
%
\section{对称函数、Young 表与组合恒等式}
\label{b5:sec:A02}

第6章用 Young 表构造 Specht 模，第7章用对称函数计算它们的特征标。这里补充表格本身的计数解释，并证明置换与成对标准表之间的对应。行表类是按行内置换得到的等价类，列交错元则属于以行表类为基的模，这两个构造分别见\secref{b5:sec:0602}和\secref{b5:sec:0603}；插入算法从带有确定行、列位置的表出发，通过移动其中的数字建立这一计数对应。

\subsection{半标准表与 Schur 多项式}
\label{b5:subsec:A02:01}

\begin{definition}\label{b5:def:semistandard-tableau}
	设 \(\lambda\) 为分拆，\(N\ge1\) 为整数，\(T\) 为在 Young 图 \([\lambda]\) 各格填入 \(1,\ldots,N\) 得到的表。若每行从左到右弱递增、每列从上到下严格递增，则称 \(T\) 为\term[banbiaozhun-Young-biao]{半标准 Young 表}[semistandard Young tableau]。若数字 \(j\) 出现 \(a_j\) 次，称 \((a_1,\ldots,a_N)\) 为其重量，并对形式变量 \(x_1,\ldots,x_N\) 记 \(x^T=\prod_{j=1}^Nx_j^{a_j}\)。
\end{definition}
\symidx{\(x^T\)：半标准 Young 表的权重单项式}

\begin{proposition}\label{b5:prop:schur-tableau-expansion}
	设 \(\lambda\) 为分拆，\(N\ge1\) 为整数，\(x_1,\ldots,x_N\) 为形式变量，\(s_\lambda\in\mathbb Q[x_1,\ldots,x_N]\) 为 Schur 多项式。则
	\[
		s_\lambda(x_1,\ldots,x_N)=\sum_T x^T,
	\]
	其中 \(T\) 遍历用 \(1,\ldots,N\) 填写的形状 \(\lambda\) 的半标准表，\(x^T=\prod_{j=1}^Nx_j^{a_j(T)}\)，\(a_j(T)\) 为数字 \(j\) 在 \(T\) 中出现的次数。
\end{proposition}
\begin{proof}
	对半标准表 \(T\)，数字不超过 \(j\) 的格子组成一个 Young 图 \(\lambda^{(j)}\)：行弱递增保证每行是前缀，列严格递增保证上面的格子先出现。同一个数字不能出现在同列，所以 \(\lambda^{(j)}/\lambda^{(j-1)}\) 为横条。反过来，由一条
	\[
		\varnothing=\lambda^{(0)}\subseteq\lambda^{(1)} \subseteq\cdots\subseteq\lambda^{(N)}=\lambda
	\]
	且每次增加横条的链，把第 \(j\) 步新增格子填为 \(j\)，就得到半标准表。这是互逆对应。

	由\cref{b5:lem:schur-cauchy}，对另一组变量 \(z\) 有
	\[
		\prod_{i,j}(1-z_ix_j)^{-1} =\sum_\lambda s_\lambda(z)s_\lambda(x).
	\]
	左端也等于 \(\prod_{j=1}^N\sum_{r\ge0}h_r(z)x_j^r\)。 使用\cref{b5:thm:pieri}逐次展开各 \(h_r(z)\) 的乘积，\(s_\lambda(z)\) 的系数恰为上述横条链的重量和。Schur 函数在每个次数构成一组基，比较系数就得到公式。
\end{proof}

例如形状 \((2,1)\) 使用 \(1,2\) 时只有两张半标准表，重量分别为 \((2,1)\)、\((1,2)\)，所以 \(s_{(2,1)}(x_1,x_2)=x_1^2x_2+x_1x_2^2\)。当允许 \(N\) 个数字时，代入 \(x_i=1\) 就计数全部半标准表；由\cref{b5:prop:schur-functor-character}，这一数目也等于 \(\mathbb S_\lambda(\mathbb C^N)\) 的维数。更细地，对角矩阵在重量 \((a_1,\ldots,a_N)\) 的张量上作用为标量 \(x_1^{a_1}\cdots x_N^{a_N}\)，而 Schur 模继承张量幂的重量分解。因此，特征标中这个单项式的系数就是相应权空间的维数，半标准表按重量计数记录了这些维数。

\begin{proposition}\label{b5:prop:A-schur-21-weights}
	设 \(V=\mathbb C^3\)，\(\mathbb S_{(2,1)}(V)\) 为相应的 Schur 模，以 \(\varepsilon_i\) 记第 \(i\) 个坐标重量。它在对角矩阵上的特征标为
	\[
		s_{(2,1)}(x_1,x_2,x_3) =\sum_{i\ne j}x_i^2x_j+2x_1x_2x_3.
	\]
	因而其维数为八，重量 \((1,1,1)\) 的权空间维数为二，其他六个重量 \(2\varepsilon_i+\varepsilon_j\)、\(i\ne j\) 的权空间均为一维。在 \(\mathfrak{sl}_3(\mathbb C)\) 中取对角 Cartan 子代数和上三角正根，基本权记为 \(\omega_1,\omega_2\)。限制到这一 Lie 代数后，它是最高权 \(\omega_1+\omega_2\) 的伴随模，重量 \((1,1,1)\) 对应其中的二维零权空间。
\end{proposition}
\begin{proof}
	形状 \((2,1)\) 的一张表可写成上行 \(a,b\)、下行 \(c\)，条件为 \(a\le b\)、\(a<c\)。可以有两格相同，但条件 \(a<c\) 排除了三格全同。若对应的重量单项式为 \(x_i^2x_j\)，\(i\ne j\)，当 \(i<j\) 时唯一的表为上行 \(i,i\)、下行 \(j\)；当 \(j<i\) 时唯一的表为上行 \(j,i\)、下行 \(i\)。这给出六张表。重量 \((1,1,1)\) 则有两张：
	\[
		\begin{array}{cc}1&2\\3&\end{array},
		\qquad
		\begin{array}{cc}1&3\\2&\end{array}.
	\]
	由表格展开式及\cref{b5:prop:schur-functor-character}，得到所述特征标与权空间维数。

	对分拆 \((2,1,0)\)，最高权由相邻行长差给出，即 \((\nu_1-\nu_2)\omega_1+(\nu_2-\nu_3)\omega_2\)。这里两差为 \(1,1\)，所以\cref{b5:thm:schur-highest-identification}将其限制识别为 \(L(\omega_1+\omega_2)\)。在 \(\mathfrak{sl}_3\) 的伴随作用中，\(E_{13}\) 是最高权为 \(\omega_1+\omega_2\) 的向量；伴随子模就是 Lie 理想，\(\mathfrak{sl}_3\) 的单性使伴随模简单，故它正是这个最高权模。对角 Cartan 子代数中的元素为 \(\operatorname{diag}(h_1,h_2,h_3)\)，其中 \(h_1+h_2+h_3=0\)，所以重量 \((1,1,1)\) 在其上取值为零。其余重量 \(2\varepsilon_i+\varepsilon_j\) 在此限制下等于 \(\varepsilon_i-\varepsilon_k\)，其中 \(k\) 为不同于 \(i,j\) 的指标，恰是六个根。于是六个一维根空间与二维 Cartan 子代数一起给出上述八维分解。
\end{proof}

\subsection{行插入与成对标准表}
\label{b5:subsec:A02:02}

对数字互异、行列均严格递增的表，插入一个未出现的数字 \(a\) 的规则如下。先在第一行找到第一个严格大于 \(a\) 的数字 \(b\)，用 \(a\) 替换它，并把 \(b\) 插入下一行；若本行没有大于待插数字的项，就把该数字附到本行末尾并停止。若下一行原不存在，就新建一行。

\begin{lemma}\label{b5:lem:row-insertion}
	设 \(T\) 为 Young 图上的填数表，所有数字互异且行列均严格递增，\(a\) 为未出现在 \(T\) 中的数字。把 \(a\) 插入第一行：每行用待插数字替换该行第一个严格大于它的数字，并将被替换的数字继续插入下一行；若没有更大的数字，就在该行末尾追加并停止，必要时新建一行。这个行插入过程保持行列严格递增，并且只在一个外角增加一格。给定插入后的表及该新增外角，可以唯一恢复原表和最初插入的数字。
\end{lemma}
\begin{proof}
	每行的替换位置取第一个大于待插数字的位置，故替换后该行仍严格递增。设上一行在第 \(c\) 列弹出了数字 \(b\)。下一行原第 \(c\) 列若存在，其数字大于 \(b\)，所以这次替换或追加的位置 \(c'\le c\)。若 \(c'<c\)，上方格子仍是上一行原第 \(c'\) 列的数字，由该行严格递增，它小于原第 \(c\) 列的 \(b\)；若 \(c'=c\)，上方的 \(b\) 已被更小的数字替换。因此新放入的 \(b\) 大于其上方数字。替换时，\(b\) 小于被替换的旧数字，而旧数字又由原表的列严格递增性小于下方原有数字，所以 \(b\) 仍小于下方数字。追加时没有下方格子：上一行长度至少达到追加位置，而本行原先比该位置短，下方各行更短。故新图仍为 Young 图且新增格为外角。

	逆过程从指定外角取出数字 \(b\)，删掉该格；在上一行找到最右边严格小于 \(b\) 的数字 \(a\)，用 \(b\) 替换它，再将 \(a\) 向上一行继续传递，直至离开第一行。所需位置存在，因为当前列的上方数字小于正在向上移动的数字。正向过程中，被插数字左边的数都小于它，右边的数都大于被弹数字；所以逆向选择的“最右边小于”恰找回原替换位置。逐行反转即恢复每一步，证明存在性及唯一性。
\end{proof}

下面的 Robinson--Schensted 对应\footnote{鲁宾逊（Gilbert de Beauregard Robinson，1906-1992），加拿大数学家，研究对称群表示和组合学；申斯特德（Craige Schensted，1927-2021），美国数学家、物理学家，提出行插入算法。}把置换配上一对表。
\begin{theorem}[Robinson--Schensted 对应]\label{b5:thm:rs-correspondence}
	设 \(n\) 为非负整数；约定 \(S_0=\{1\}\)，空分拆的 Young 图有唯一的空标准表。置换 \(w\in S_n\) 与一对同形状、各用 \(1,\ldots,n\) 填数的标准 Young 表 \((P,Q)\) 一一对应：按 \(w_i=w(i)\) 的次序行插入，得到插入表 \(P\)，并在每步的新增格记录步数，得到记录表 \(Q\)。当 \(n=0\) 时，唯一的置换对应空表对 \((\varnothing,\varnothing)\)。
\end{theorem}
\begin{proof}
	当 \(n=0\) 时，两边都只有一个对象。以下设 \(n\ge1\)。依次把 \(w_1,\ldots,w_n\) 行插入到空表，得到 \(P\)。第 \(i\) 步新增一格，在第二张图的相同位置写上 \(i\)，得到记录表 \(Q\)。前一引理保证 \(P\) 为标准表；各步形状都是 Young 图，格子右边或下边的格子必后加入，所以 \(Q\) 也是标准表。

	反过来，从 \(Q\) 中找到最大数字 \(n\) 所在的外角，在 \(P\) 中从同一外角作逆插入，得到数字 \(w_n\)，并删去 \(Q\) 的这格。递减重复，得到 \(w_{n-1},\ldots,w_1\)。每个数字恰从原来填有 \(1,\ldots,n\) 的 \(P\) 中移出一次，所以输出是一个置换。引理的逐步互逆性说明两个整体过程互逆。
\end{proof}

\Cref{b5:alg:robinson-schensted}把这两个互逆过程写成可逐格执行的步骤。正向新增的格由记录表保存，逆向按记录数字从大到小删除外角。

\begin{algorithm}[htbp]
	\caption{Robinson--Schensted 行插入与逆插入}
	\label{b5:alg:robinson-schensted}
	\begin{algorithmsteps}{置换 \(w\in S_n\)，或一对同形状、各填 \(1,\ldots,n\) 的标准 Young 表 \((P,Q)\)。}{输入为置换时输出对应的表对 \((P,Q)\)；输入为表对时输出置换 \(w\)。}
		\If{输入为置换 \(w\)}
			\State 置 \(P,Q\) 为空表。
			\For{\(i=1,\ldots,n\)}
				\State 从首行起把 \(w_i\) 插入 \(P\)：在每行替换第一个更大的数，把被替换数送往下一行；若没有更大的数，就在该行末添格并停止。记新增格为 \(c_i\)。
				\State 在 \(Q\) 的同一格 \(c_i\) 填入 \(i\)。
			\EndFor
			\State 输出 \((P,Q)\)。
		\Else
			\For{\(i=n,n-1,\ldots,1\)}
				\State 找到 \(Q\) 中数字 \(i\) 所在的外角，记其行为 \(r\)，并删去该格；从 \(P\) 的同一格取出数字 \(b\) 并删格。
				\State 依次经过第 \(r-1,\ldots,1\) 行：找本行最右边小于 \(b\) 的数 \(a\)，用 \(b\) 替换它，再令 \(b\leftarrow a\)。最后置 \(w_i=b\)。
			\EndFor
			\State 输出 \(w\)。
		\EndIf
	\end{algorithmsteps}
\end{algorithm}

每次正向插入至多新建一行，逆向每步都向上一行移动，因此两种操作均在有限步终止；其互逆性由前述引理保证。

例如 \(w=312\) 的插入过程先得到单格 \(3\)，再把 \(1\) 插入并弹出 \(3\)，最后将 \(2\) 附到第一行。于是
\[
	P=\begin{array}{cc}1&2\\3&\end{array},
	\qquad
	Q=\begin{array}{cc}1&3\\2&\end{array}.
\]
从 \(Q\) 的格子 \(3,2,1\) 依次逆插入，正好恢复 \(2,1,3\) 的逆读取次序。

\begin{proposition}\label{b5:prop:A-rs-2413}
	对排列 \(w=2413\in S_4\)，按行插入约定得到
	\[
		P(w)=\begin{array}{cc}1&3\\2&4\end{array},
		\qquad
		Q(w)=\begin{array}{cc}1&2\\3&4\end{array}.
	\]
	它的逆排列为 \(w^{-1}=3142\)，且
	\[
		P(w^{-1})=Q(w),\qquad Q(w^{-1})=P(w).
	\]
	这两个排列的最长严格递增子序列和最长严格递减子序列的长度都为二，与公共形状 \((2,2)\) 的首行、首列长度相等。
\end{proposition}
\begin{proof}
	对 \(2413\)，前两步得到首行 \(2,4\)。插入 \(1\) 后，首行成为 \(1,4\)，弹出的 \(2\) 组成第二行；插入 \(3\) 时，首行的 \(4\) 被弹到第二行末尾，得到所述 \(P(w)\)。四次新增格依次为 \((1,1),(1,2),(2,1),(2,2)\)，故 \(Q(w)\) 如上。

	逆插入时，从 \(Q\) 的格子 \(4\) 所在位置取出 \(P\) 中的 \(4\)，用它替换首行最右边小于四的 \(3\)，输出 \(w_4=3\)。接着从格子 \(3\) 取出 \(2\)，用它替换首行的 \(1\)，输出 \(w_3=1\)。余下首行为 \(2,4\)，再依次输出 \(w_2=4\)、\(w_1=2\)，恢复原排列。

	交换位置与数值得到逆排列 \(3142\)。它的插入表依次为
	\[
		\begin{array}{c}3\end{array},\qquad
		\begin{array}{c}1\\3\end{array},\qquad
		\begin{array}{cc}1&4\\3&\end{array},\qquad
		\begin{array}{cc}1&2\\3&4\end{array}.
	\]
	新增格依次为 \((1,1),(2,1),(1,2),(2,2)\)，得到 \(Q(w^{-1})=P(w)\)。在 \(2413\) 中，\(2,4\) 和 \(4,1\) 分别是长度二的递增、递减子序列；四种三项子序列 \(241,243,213,413\) 均不严格单调。在 \(3142\) 中，\(1,4\) 和 \(3,1\) 给出长度二的两类子序列，而三项子序列 \(314,312,342,142\) 也均不严格单调。因此两类最大长度都是二。
\end{proof}

首行采用“替换第一个更大数字”的规则，是为了保存各长度递增子序列可能达到的最小末项。较小的末项更容易接上以后出现的数字；替换它只改善延长的可能性，不必保存整个子序列。

\begin{proposition}\label{b5:prop:A-schensted-first-row}
	设 \(w=w_1\cdots w_n\in S_n\)，其 Robinson--Schensted 插入表的形状为 \(\lambda\)。则 \(\lambda_1\) 等于 \(w\) 的最长严格递增子序列的长度；空图的首行长度与空排列的子序列长度均规定为零。
\end{proposition}
\begin{proof}
	读取任意前缀后，把插入表首行写为 \(b_1<\cdots<b_r\)。归纳证明：\(b_j\) 恰是该前缀中所有长度为 \(j\) 的严格递增子序列的最小末项，且没有更长的递增子序列。空前缀满足此断言。

	新读入 \(a\) 时，设 \(j\) 为第一个满足 \(b_j>a\) 的位置；若不存在，令 \(j=r+1\)。当 \(j>1\) 时，长度 \(j-1\) 的某条递增子序列以 \(b_{j-1}<a\) 结尾，可以接上 \(a\)；当 \(j=1\) 时，单项 \(a\) 本身即可。故 \(a\) 成为长度 \(j\) 的新最小末项。对 \(k<j\)，旧最小末项 \(b_k<a\)，不会被改善；对 \(k>j\)，若新出现长度 \(k\) 的递增子序列，它必须以 \(a\) 结尾，其前 \(k-1\) 项的末项却至少为 \(b_{k-1}>a\)，矛盾。没有第 \(k-1\) 个旧末项时，这样的前缀子序列更不存在。因此，首行只需在位置 \(j\) 替换为 \(a\)，或在 \(j=r+1\) 时追加 \(a\)，归纳条件保持。读完全部输入即得结论。
\end{proof}

Schensted 定理还断言，公共形状 \(\lambda\) 的首列长度 \(\lambda'_1\) 等于最长严格递减子序列的长度\cite[Theorems 1--2, pp. 183, 188]{Schensted1961Subsequences}。递减情形用到另一项插入性质：把输入次序倒转，所得插入表是原插入表的转置；倒转又把递减子序列变成递增子序列，所以可归结为首行长度的计算。一般的行、列插入交换论证见同文\cite[Lemmas 6--7, pp. 183--188]{Schensted1961Subsequences}。

\ProseParagraphBreak
取逆交换两张表也对任意 \(w\in S_n\) 成立：
\[
	P(w^{-1})=Q(w),\qquad Q(w^{-1})=P(w).
\]
它的一般证明把输入改写为置换矩阵：矩阵转置对应排列取逆，而相应图构造的对称性使两张表交换\cite[Section 4, Theorem 3, p. 719]{Knuth1970Tableaux}。这一步需要该文建立的图构造，逐次插入 \(2413\) 的计算只是这一对称性的一个实例。由它及已证的双射，\(w=w^{-1}\) 当且仅当 \(P(w)=Q(w)\)：一个方向是上述对称性，另一个方向则由同一表对只能对应一个置换得到。记 \(f^\lambda\) 为形状 \(\lambda\) 的标准表数，则 \(S_n\) 中满足 \(w^2=1\) 的置换总数为 \(\sum_{\lambda\vdash n}f^\lambda\)，因为每张标准表 \(T\) 恰给出一个表对 \((T,T)\)。这里计入恒等置换；若只计阶恰为二的对合，应将这个总数减一。

\subsection{维数恒等式与格路径计数}
\label{b5:subsec:A02:03}

\begin{corollary}\label{b5:cor:tableau-squares}
	设 \(n\) 为非负整数，对每个分拆 \(\lambda\vdash n\)，记 \(f^\lambda\) 为形状 \(\lambda\) 的标准表数，则
	\[
		\sum_{\lambda\vdash n}(f^\lambda)^2=n!.
	\]
\end{corollary}
\begin{proof}
	对固定形状，可独立选择 \(P,Q\)，共有 \((f^\lambda)^2\) 对。上一对应把所有这些对恰好与 \(n!\) 个置换对应，求和即得。
\end{proof}

群代数 \(\mathbb C[S_n]\) 的半单分解也给出这一维数恒等式；插入法则直接给出组合证明。

\ProseParagraphBreak
两行形状 \((a,b)\)、\(a\ge b\ge1\) 还有直接的路径计数。按 \(1,\ldots,a+b\) 的次序读取标准表，把数字位于第一行记为一步向上，位于第二行记为一步向下，路径从高度零出发。设两行的数字分别为 \(a_1<\cdots<a_a\)、\(b_1<\cdots<b_b\)。列严格递增要求 \(a_i<b_i\)；读到 \(b_i\) 时，这保证至少已读到 \(i\) 个第一行数字。反过来，若每个前缀的向上步数不少于向下步数，读到 \(b_i\) 时第一行至少已有 \(i\) 个数，且数字互异，故 \(a_i<b_i\)。这就证明了两条件等价。总路径数为 \(\binom{a+b}{b}\)。对第一次越到高度 \(-1\) 的坏路径，把到该点为止的上、下步互换，就得到从高度零出发、具有 \(a+1\) 个上步和 \(b-1\) 个下步的路径，其终点高度为 \(a-b+2\ge2\)。反过来，每条这样的路径必经过高度 \(1\)；把到第一次达到高度 \(1\) 为止的上、下步互换，便得到第一次在该位置越到高度 \(-1\) 的坏路径。这两次反射互逆，因此
\[
	f^{(a,b)}=\binom{a+b}{b}-\binom{a+b}{b-1}.
\]
当 \(a=b=r\) 时，这是 Catalan 数
\[
	f^{(r,r)}=\frac1{r+1}\binom{2r}{r}.
\]
该推导与\cref{b5:thm:hook-length}给出的两行钩长计算一致，但它还解释了相减的那项在数什么。

\ProseParagraphBreak
例如形状 \((3,3)\) 的坏路径 \(U,D,D,U,U,D\)，其中 \(U\) 表示上步、\(D\) 表示下步，前三步的高度依次为 \(1,0,-1\)。反射这个首次越界的前缀 \(U,D,D\)，就得到 \(D,U,U,U,U,D\)，有四个上步和两个下步；反过来，它第一次达到高度一也在第三步，反射前三步便恢复原坏路径。因此坏路径由无约束的四上二下路径计数，共有 \(\binom62=15\) 条；全部三上三下路径有 \(\binom63=20\) 条，合法路径与标准表就各有 \(20-15=5\) 个。

% !TeX root = ../../Book5.tex
% 纲要标题：### A.3 代数组合中的计数方法
% 纲要位置：工作记录/计划纲要.md，第 4933 行。
% 纲要细目（保留原定写作范围）：
%
% #### A.3.1 轨道计数与带权循环指标
% #### A.3.2 分拆与逆序数的生成函数
% #### A.3.3 有限域子空间与 Gaussian 二项式
%
\section{代数组合中的计数方法}
\label{b5:sec:A03}

用表示或生成函数描述一个集合后，计数往往可以通过取迹、取系数或比较基来完成。使用这些方法时，须说明每个代数项所代表的对象。下面分别讨论对称性、长度和有限域上的子空间。

\subsection{轨道计数与带权循环指标}
\label{b5:subsec:A03:01}

\begin{proposition}\label{b5:prop:weighted-orbit-count}
	设有限群 \(G\) 作用于有限集 \(X\)，记 \(X/G\) 为轨道集、\(X^g=\{x\in X:gx=x\}\) 为 \(g\in G\) 的不动点集，则
	\[
		|X/G|=\frac1{|G|}\sum_{g\in G}|X^g|.
	\]
	若每个 \(x\in X\) 另带一个在轨道上恒定的有理数系数多项式环 \(\mathbb Q[z_1,\ldots,z_r]\) 中的单项式权重 \(m(x)\)，记其在轨道 \(O\) 上的共同值为 \(m(O)\)，则轨道的权重和为
	\[
		\sum_{O\in X/G}m(O) =\frac1{|G|}\sum_{g\in G}\sum_{x\in X^g}m(x).
	\]
\end{proposition}
\begin{proof}
	固定一个轨道 \(O\)。右端双和中该轨道的贡献是
	\[
		\frac1{|G|}\sum_{x\in O}|G_x|m(O).
	\]
	轨道稳定子群公式使每项 \(|G_x|=|G|/|O|\)，故贡献恰为 \(m(O)\)。逐个轨道相加就得到带权式；令全部权重为一得到第一式。
\end{proof}

设 \(G\le S_n\) 置换 \(n\) 个位置，用 \(r\) 种颜色染色，并用 \(z_i\) 记录颜色 \(i\) 出现次数。若 \(g\) 有 \(m_j(g)\) 个长度 \(j\) 的轮换，固定染色必须在每个轮换上同色，因此轨道重量生成函数为
\begin{equation}\label{b5:eq:polya-cycle-index}
	\frac1{|G|}\sum_{g\in G} \prod_{j\ge1}(z_1^j+\cdots+z_r^j)^{m_j(g)}.
\end{equation}
这就是\term[xunhuan-zhibiao]{循环指标}[cycle index]在染色问题中的替换规则，它不要求各个染色轨道的大小相同：带权公式对每个轨道只计入一次共同权重。

\ProseParagraphBreak
例如四个位置只允许旋转、使用红蓝两色时，四个旋转的循环型为 \(1^4,4,2^2,4\)。所以计数多项式为
\[
	\frac14\bigl((a+b)^4+(a^2+b^2)^2+2(a^4+b^4)\bigr) =a^4+a^3b+2a^2b^2+ab^3+b^4.
\]
共有六类，两红两蓝有两类，分别是同色相邻与红蓝交替。若还允许翻转或全体颜色交换，作用群就改变了，必须重新计算。

\subsection{分拆与逆序数的生成函数}
\label{b5:subsec:A03:02}

整数分拆可以按每个部分的出现次数计数。每个正整数 \(j\) 可出现 \(0,1,\ldots\) 次，所以形式级数恒等式为
\[
	\sum_{n\ge0}p(n)t^n=\prod_{j\ge1}(1-t^j)^{-1}.
\]
固定 \(t^n\) 时只有 \(j\le n\) 的因子有贡献，故右端在形式意义下良定。将只出现零或一次的选择编码为 \(1+t^j\)，又得到
\[
	\prod_{j\ge1}(1+t^j) =\prod_{j\ge1}\frac{1-t^{2j}}{1-t^j} =\prod_{\substack{j\ge1\\j\ \mathrm{odd}}}(1-t^j)^{-1}.
\]
因此不同部分的分拆与全部部分为奇数的分拆等数。这是逐次形式消去，不能解读为任意数值 \(t\) 下都可重排的无穷乘积。

\ProseParagraphBreak
这个等数还有一个显式双射。对只含奇数部分的分拆，将每个奇数 \(d\) 的重数写成二进制展开
\[
	m_d=\sum_{r\ge0}\varepsilon_{d,r}2^r, \qquad \varepsilon_{d,r}\in\{0,1\}.
\]
每当 \(\varepsilon_{d,r}=1\)，就将 \(2^r\) 个 \(d\) 合并为一个部分 \(2^rd\)。不同的 \((d,r)\) 不会产生相同部分，因为每个正整数都唯一分解为二的幂乘奇数；故所得分拆的部分互不相同。反过来，把每个不同部分 \(2^rd\) 拆成 \(2^r\) 个 \(d\)，便恢复原来的奇数分拆。两种操作都保持总和，而且互为逆。例如
\[
	3+1+1+1\ \longleftrightarrow\ 3+2+1, \qquad 6+4+1\ \longleftrightarrow\ 3+3+1+1+1+1+1.
\]
生成函数中的同一机制是
\[
	\prod_{r\ge0}(1+t^{2^rd})=(1-t^d)^{-1} \qquad(d\;\text{为正奇数}).
\]
左端选择重数的各个二进制位，右端直接记录任意非负重数。固定次数时只涉及有限多位，所以这一解释也符合形式级数的运算规则。

\begin{proposition}\label{b5:prop:permutation-length-enumerator}
	设 \(n\) 为非负整数，\(q\) 为形式变量，\(\ell(w)\) 为 \(w\in S_n\) 关于邻位换位的长度。则在 \(\mathbb Z[q]\) 中有
	\[
		\sum_{w\in S_n}q^{\ell(w)} =\prod_{j=1}^n(1+q+\cdots+q^{j-1}).
	\]
\end{proposition}
\begin{proof}
	由\cref{b5:prop:permutation-length-inversions}，长度等于逆序数。在一个 \(S_{n-1}\) 置换的单行记号中，任取一个空隙放入最大数字 \(n\)；这里不使用上一节的 Young 表行插入。若它右边有 \(r\) 个数字，就新增恰 \(r\) 个逆序，而 \(r\) 唯一遍历 \(0,\ldots,n-1\)。因此长度多项式递推为前一级乘 \(1+q+\cdots+q^{n-1}\)。从 \(n=1\) 的初值一开始便得到公式。
\end{proof}

在 \(q=1\) 处恢复 \(n!\)，在 \(q=-1\) 处（\(n\ge2\)）得到零，表示奇、偶置换等数。最高次数为 \(n(n-1)/2\)，最高项系数一，正对应唯一的最长置换。这个多项式记录了各长度的置换数，后面还会出现在旗空间的计数中。

\subsection{有限域子空间与 Gaussian 二项式}
\label{b5:subsec:A03:03}

\begin{proposition}\label{b5:prop:gaussian-subspace-count}
	设 \(n,r\) 为满足 \(0\le r\le n\) 的整数，\(Q\) 为形式变量。在 \(\mathbb Q(Q)\) 中定义
	\[
		{n\brack r}_Q=\prod_{i=1}^r\frac{1-Q^{n-r+i}}{1-Q^i},
	\]
	取空乘积为一。则 \({n\brack r}_Q\in\mathbb Z[Q]\)。对任意素数幂 \(q\)，记 \(\mathbb F_q\) 为 \(q\) 元有限域，\({n\brack r}_q\) 为此多项式在 \(Q=q\) 处的值；它等于 \(\mathbb F_q^n\) 中 \(r\) 维子空间的数目，并且
	\[
		{n\brack r}_q =\frac{(q^n-1)(q^n-q)\cdots(q^n-q^{r-1})} {(q^r-1)(q^r-q)\cdots(q^r-q^{r-1})}.
	\]
\end{proposition}
\symidx{\({n\brack r}_Q\)：Gaussian 二项式}
\begin{proof}
	先固定素数幂 \(q\)，将 \(\mathbb F_q^n\) 中 \(r\) 维子空间的数目记为 \(N_{n,r}(q)\)。有序线性无关 \(r\) 元组可逐个选择：第 \(i+1\) 个向量不得落在前 \(i\) 个生成的 \(q^i\) 元子空间中，所以共有上述分子的数目。每个固定 \(r\) 维子空间的有序基数为分母，且每个独立元组唯一确定其生成子空间。相除后，提出相应的 \(q^i\) 并调整指标，得到
	\[
		N_{n,r}(q)=\prod_{i=1}^r\frac{1-q^{n-r+i}}{1-q^i}.
	\]
	当 \(1\le r\le n-1\) 时，固定超平面 \(H=\mathbb F_q^{n-1}\subset V=\mathbb F_q^n\)，将 \(r\) 维子空间分成包含于 \(H\) 和不包含于 \(H\) 两类。后一类的交 \(U=W\cap H\) 维数为 \(r-1\)。对每个固定的 \(U\)，\(W/U\) 是 \(V/U\) 中不落在超平面 \(H/U\) 中的一条直线；先在 \(H/U\) 中选基，再添一个向量得到 \(V/U\) 的基，使 \(H/U\) 恰由最后一坐标为零刻画。将直线代表的最后一坐标归一为一后，其余 \(n-r\) 个坐标可任取，故这样的直线共有 \(q^{n-r}\) 条。于是
	\[
		N_{n,r}(q)=N_{n-1,r}(q)+q^{n-r}N_{n-1,r-1}(q).
	\]
	零子空间和整个空间各有一个，故 \(N_{n,0}(q)=N_{n,n}(q)=1\)。用相同的边界值 \(F_{n,0}(Q)=F_{n,n}(Q)=1\) 及递推
	\[
		F_{n,r}(Q)=F_{n-1,r}(Q)+Q^{n-r}F_{n-1,r-1}(Q) \qquad(1\le r\le n-1)
	\]
	逐级定义 \(F_{n,r}(Q)\)。归纳得到 \(F_{n,r}(Q)\in\mathbb Z[Q]\)，且对每个素数幂 \(q\) 都有 \(F_{n,r}(q)=N_{n,r}(q)\)。将有理函数 \({n\brack r}_Q-F_{n,r}(Q)\) 通分，其分子是一个有理系数多项式。分母在所有素数幂 \(q>1\) 处均不为零，而分子在这些无限多个值处均为零，故分子恒为零。这就证明 \({n\brack r}_Q=F_{n,r}(Q)\)，以及形式多项式递推
	\begin{equation}\label{b5:eq:gaussian-recursion}
		{n\brack r}_Q={n-1\brack r}_Q +Q^{n-r}{n-1\brack r-1}_Q \qquad(1\le r\le n-1).
	\end{equation}
\end{proof}

在多项式中代入 \(Q=1\)，递推变成普通二项式递推，边界值也相同，因此
\[
	\left.{n\brack r}_Q\right|_{Q=1}=\binom nr.
\]
这里是在形式多项式中求值，没有引入单元素的域。

\begin{proposition}\label{b5:prop:gaussian-young-area}
	设 \(0\le r\le n\) 为整数，\(Q\) 为形式变量。以 \(\lambda\subseteq r\times(n-r)\) 表示 Young 图至多有 \(r\) 行、每行至多有 \(n-r\) 个格子的分拆，包括空分拆。则
	\[
		{n\brack r}_Q =\sum_{\lambda\subseteq r\times(n-r)}Q^{|\lambda|}.
	\]
	因此 \(Q^d\) 的系数是上述矩形内恰有 \(d\) 个格子的 Young 图数目。
\end{proposition}
\begin{proof}
	将右端记为 \(G_{n,r}(Q)\)。当 \(r=0\) 或 \(r=n\) 时，只有空分拆，所以边界值是一。对 \(1\le r\le n-1\)，按第一行是否达到宽度 \(n-r\) 分类。没有达到的 Young 图位于 \(r\times(n-r-1)\) 的矩形内，贡献 \(G_{n-1,r}(Q)\)。达到的图删去第一行后，留下 \((r-1)\times(n-r)\) 内的任意 Young 图，而删去的 \(n-r\) 个格子贡献因子 \(Q^{n-r}\)。故
	\[
		G_{n,r}(Q)=G_{n-1,r}(Q)+Q^{n-r}G_{n-1,r-1}(Q).
	\]
	这与\eqref{b5:eq:gaussian-recursion}的递推及边界值相同，归纳即得等式。
\end{proof}

例如在 \(2\times2\) 矩形内只有
\[
	\varnothing,\quad(1),\quad(2),\quad(1,1),\quad(2,1),\quad(2,2),
\]
格子数依次为 \(0,1,2,2,3,4\)，所以
\[
	{4\brack2}_Q=1+Q+2Q^2+Q^3+Q^4, \qquad {4\brack2}_2=35.
\]
次数二的系数为二，是因为 \((2)\) 与 \((1,1)\) 两种图具有相同面积。

\ProseParagraphBreak
Young 图面积还可以直接数出有限域子空间的自由坐标。把一个 \(r\) 维子空间的基作为行向量，化为唯一的最简行阶梯矩阵，设主元列为 \(p_1<\cdots<p_r\)。第 \(i\) 行主元右边共有 \(n-p_i\) 列，其中 \(r-i\) 列是别行的主元列，必须为零；其余
\[
	\lambda_i=n-r+i-p_i
\]
个位置可任意取值。这些数满足 \(n-r\ge\lambda_1\ge\cdots\ge\lambda_r\ge0\)，因而给出矩形内的一个 Young 图。反过来，\(p_i=n-r+i-\lambda_i\) 恢复主元列，所以两种资料一一对应。对固定图，全部自由位置数是 \(|\lambda|\)，在 \(\mathbb F_q\) 上给出 \(q^{|\lambda|}\) 个子空间。由此能逐项解释 Gaussian 二项式的有限域求值，而不只是比较总数。

\ProseParagraphBreak
完整旗 \(0=V_0\subset V_1\subset\cdots\subset V_n=\mathbb F_q^n\)、 \(\dim V_i=i\)，可先选一条直线，再在商空间中选完整旗。因此旗的数目为
\[
	\prod_{j=1}^n\frac{q^j-1}{q-1} =\prod_{j=1}^n(1+q+\cdots+q^{j-1}).
\]
它与置换长度多项式相同，提示了长度与旗几何之间的联系；两边在此均有独立的计数证明。

\begin{proposition}\label{b5:prop:flag-permutation-coordinates}
	设 \(q\) 为素数幂，\(n\ge1\)，\(B\le\operatorname{GL}_n(\mathbb F_q)\) 为可逆上三角矩阵组成的子群。\(\mathbb F_q^n\) 中的完整旗按 \(w\in S_n\) 分为 \(B\) 的轨道 \(C_w\)，其中 \(C_w\) 包含置换旗
	\[
		V_j=\langle e_{w(1)},\ldots,e_{w(j)}\rangle.
	\]
	每个 \(C_w\) 中的旗有唯一的一组规范代表向量
	\[
		u_j=e_{w(j)}+ \sum_{\substack{i<w(j)\\ i\notin\{w(1),\ldots,w(j-1)\}}} a_{ij}e_i, \qquad V_j=\langle u_1,\ldots,u_j\rangle.
	\]
	各系数 \(a_{ij}\in\mathbb F_q\) 独立任取，故 \(|C_w|=q^{\ell(w)}\)。
\end{proposition}
\begin{proof}
	依次在 \(V_j\) 中选取向量，并消去已经选定的主元坐标。前 \(j-1\) 个代表在这些坐标上的矩阵是对角元为一的三角矩阵，因而可逆；所以 \(V_j\) 中这些坐标全为零的向量构成一维空间。取这个空间中最后一个非零坐标的位置为 \(w(j)\)，并将其值归一为一，就唯一得到 \(u_j\)。该位置与之前的主元不同，其右边坐标全为零，于是恰有陈述中的形式。\(n\) 个不同主元组成一个置换 \(w\)；反过来，任意给定的这些系数都产生一组基和一个完整旗。

	为看出它们恰是 \(B\) 的轨道，将 \(e_{w(j)}\) 送到 \(u_j\) 所得矩阵是对角元为一的上三角矩阵，故全部这样的旗都在置换旗的 \(B\) 轨道中。上三角变换保持标准子空间 \(E_i=\langle e_1,\ldots,e_i\rangle\)，而这些旗满足
	\[
		\dim(V_j\cap E_i)=\#\{\,k\le j\mid w(k)\le i\,\}.
	\]
	这是因为不同主元之间不能消去最后的非零坐标；取值矩阵中的这些维数也唯一恢复每个 \(w(j)\)。因此不同 \(w\) 不可能属于同一轨道。最后，第 \(j\) 个向量的自由坐标数为
	\[
		\#\{\,k>j\mid w(k)<w(j)\,\}.
	\]
	对 \(j\) 求和就是逆序数，也就是\cref{b5:prop:permutation-length-inversions}中的 \(\ell(w)\)。
\end{proof}

在 \(n=2\) 时，这就是\subsecref{b5:subsec:A01:03}中的固定直线与仿射直线轨道，分别有 \(1\) 与 \(q\) 个旗。在 \(n=3\) 时，把旗记为上述三个代表向量，六种形式为
\[
	\begin{aligned}
		123&:\ (e_1,e_2,e_3),\\
		213&:\ (e_2+a e_1,e_1,e_3),\\
		132&:\ (e_1,e_3+a e_2,e_2),\\
		231&:\ (e_2+a e_1,e_3+b e_1,e_1),\\
		312&:\ (e_3+a e_1+b e_2,e_1,e_2),\\
		321&:\ (e_3+a e_1+b e_2,e_2+c e_1,e_1).
	\end{aligned}
\]
其中 \(a,b,c\) 都在 \(\mathbb F_q\) 中独立取值，故六个轨道依次贡献 \(1,q,q,q^2,q^2,q^3\)，总数是 \(1+2q+2q^2+q^3\)。这些自由坐标说明了旗计数中每个 \(q^{\ell(w)}\) 的来源。

% !TeX root = ../../Book5.tex
% 纲要标题：### A.4 Hecke 代数与 Kazhdan--Lusztig 理论
% 纲要位置：工作记录/计划纲要.md，第 4941 行。
% 纲要细目（保留原定写作范围）：
%
% #### A.4.1 参数、二次关系与标准基
% #### A.4.2 反演算子与规范基
% #### A.4.3 胞腔与一个二维表示
%
\section{Hecke 代数与 Kazhdan--Lusztig 理论}
\label{b5:sec:A04}

群代数中简单反射满足平方为一。如果把这个二次关系变形，同时保留辫关系，就得到 Hecke 代数。它的标准基仍由群元素标号，却还存在与一个参数反演相容的特殊基。本节固定一种归一化，并通过 \(S_3\) 把不同基与胞腔表示具体算出来。

\subsection{参数、二次关系与标准基}
\label{b5:subsec:A04:01}

这个变形的二次关系可以从有限域上的旗直接算出。设 \(q\) 为素数幂，\(G=\operatorname{GL}_2(\mathbb F_q)\)，\(X\) 为 \(\mathbb F_q^2\) 中直线的集合。每条直线有 \(q-1\) 个非零向量，故 \(|X|=q+1\)。在函数空间 \(V=\mathbb C^X\) 上，\(G\) 按 \((g\cdot f)(L)=f(g^{-1}L)\) 作用，定义
\[
(Tf)(L)=\sum_{\substack{L'\in X\\L'\ne L}}f(L').
\]
群作用保持两条直线是否相同，因此 \(T\) 与 \(G\) 的作用交换。它的矩阵在对角线上为零，其余位置为一。\(T^2\) 的对角系数是在两步之后回到原直线的中间选择数，共有 \(q\) 种；非对角系数要求中间直线不同于两个端点，共有 \(q-1\) 种。因此
\[
T^2=(q-1)T+qI,\qquad (T-qI)(T+I)=0.
\]
矩阵 \(M\) 与 \(G\) 的作用交换，等价于其条目满足 \(M(gL_1,gL_2)=M(L_1,L_2)\)，即在 \(X\times X\) 的各轨道上取常数。\(G\) 对直线传递，且任意两条不同直线可选代表组成基，故能把任一有序不同直线对送到另一对。因此恰有对角与非对角两条轨道，得到
\(\operatorname{End}_G(V)=\mathbb C I\oplus\mathbb C T\)。于是这个二次式决定了整个相容算子代数的乘法。当 \(q=2\) 时，\(\operatorname{GL}_2(\mathbb F_2)\) 在三条直线上的作用将它识别为 \(S_3\)，一条直线的稳定子群为阶二子群。因此这里的 \(T=J-I\) 正是第四章\subsecref{b5:subsec:0402:02}中的双陪集算子，关系也成为 \(T^2=T+2I\)。该例的卷积测度把稳定子群的体积归一为一，与这里每条陪集贡献一次的计数一致。

\ProseParagraphBreak
取数值 \(v=\sqrt q\)，令 \(H=v^{-1}T\)，上述关系成为
\[
H^2=(v-v^{-1})H+I,\qquad (H-vI)(H+v^{-1}I)=0.
\]
有限域的大小由此给出参数 \(q=v^2\)；把 \(v\) 改为形式变量，就能在同一个系数环上研究这个乘法随参数的变化。

\begin{definition}\label{b5:def:hecke-algebra}
设 \((W,S)\) 为 Coxeter 系统，\(v\) 为形式变量，\(R=\mathbb Z[v,v^{-1}]\)。用生成元 \(H_s\)（\(s\in S\)）、关系
\[
(H_s-v)(H_s+v^{-1})=0
\]
以及与 Coxeter 系统相同的辫关系定义的含单位结合 \(R\) 代数，称为等参数\term[Hecke-daishu]{Hecke 代数}[Hecke algebra]，记为 \(\mathcal H(W)\)。
\end{definition}
\symidx{\(\mathcal H(W)\)：Coxeter 系统的等参数 Hecke 代数}

\begin{theorem}[Hecke 标准基]\label{b5:thm:hecke-standard-basis}
设 \((W,S)\) 为 Coxeter 系统，\(\ell\) 为关于 \(S\) 的长度函数，\(v\) 为形式变量，\(R=\mathbb Z[v,v^{-1}]\)。令 \(\mathcal H(W)\) 为 \(R\) 上的等参数 Hecke 代数，以 \(H_s\)（\(s\in S\)）为生成元，满足 \((H_s-v)(H_s+v^{-1})=0\) 及该 Coxeter 系统的辫关系。对约化词 \(w=s_1\cdots s_r\)，乘积
\(H_w=H_{s_1}\cdots H_{s_r}\) 与约化词的选择无关，而且 \(\{H_w\mid w\in W\}\) 为 \(R\) 基。乘法满足
\begin{equation}\label{b5:eq:hecke-standard-multiplication}
H_sH_w=
\begin{cases}
H_{sw},&\ell(sw)=\ell(w)+1,\\
H_{sw}+(v-v^{-1})H_w,&\ell(sw)=\ell(w)-1.
\end{cases}
\end{equation}
\end{theorem}
\symidx{\(H_w\)：Hecke 代数的标准基元}
\begin{proof}
记 \(d=v-v^{-1}\)。由\cref{b5:thm:coxeter-word-properties}的辫词性质，\(H_w\) 良定义。若 \(sw\) 较短，则 \(w=s(sw)\) 是约化分解，二次关系 \(H_s^2=1+dH_s\) 给出第二种乘法；较长时直接接上 \(s\)。因此所有生成元乘积都在这些 \(H_w\) 的生成空间中。

为证线性无关，要先在一个已经具有自由基的模上实现这些乘法。普通群代数的正则表示给出 \(L_se_w=e_{sw}\)、\(R_te_w=e_{wt}\)，左右乘法交换，且单位基向量 \(e_1\) 能生成全部基向量。这里尝试保留这个构造，只在长度下降时加入二次关系要求的 \(d\) 项。取自由模 \(E=\bigoplus_{w\in W}Re_w\)，按同一长度规则定义算子
\[
 L_se_w=e_{sw}+\begin{cases}0,&\ell(sw)>\ell(w),\\de_w,&\ell(sw)<\ell(w),\end{cases}
 \quad
 R_te_w=e_{wt}+\begin{cases}0,&\ell(wt)>\ell(w),\\de_w,&\ell(wt)<\ell(w).\end{cases}
\]
例如在秩一的 \(W=\{1,s\}\) 中，有 \(L_se_1=e_s\)、\(L_se_s=e_1+de_s\)，所以
\[
[L_s]_{(e_1,e_s)}=\begin{pmatrix}0&1\\1&d\end{pmatrix},
\qquad L_s^2=I+dL_s.
\]
当 \(v=1\) 时，\(d=0\)，这正是交换两个基向量的群正则表示。一般情形在 \(e_w,e_{sw}\) 上也直接算得 \(L_s^2=1+dL_s\)。右算子的用途是把辫关系在所有 \(e_w\) 上的验证归结到 \(e_1\) 上；为此先验证它们与左算子交换。

记 \(n=\ell(w)\)，先设 \(sw\ne wt\)。当 \(\ell(sw),\ell(wt)\) 分别为 \(n+1,n-1\) 或 \(n-1,n+1\) 时，\(\ell(swt)\) 与这两个数都相差一，故只能为 \(n\)。当两者都为 \(n+1\) 时，若 \(\ell(swt)=n\)，取 \(w\) 的约化词并在末尾接 \(t\)，得到 \(wt\) 的约化词。对 \((wt)^{-1}\) 右乘 \(s\) 应用交换性质，再取逆，便得到 \(swt\) 可由该词删去一位；这把右乘版本转为左乘版本。若删去 \(w\) 中的位置，再右乘 \(t\) 就得到长度至多为 \(n-1\) 的 \(sw\)，矛盾；若删去末尾的 \(t\)，则 \(swt=w\)，即 \(sw=wt\)，仍矛盾。因此此时 \(\ell(swt)=n+2\)。

当 \(\ell(sw)=\ell(wt)=n-1\) 时，若 \(\ell(swt)=n\)，令 \(z=sw\)。则 \(sz=w\)、\(zt=swt\) 都长 \(n\)，且二者不同。刚证的两次长度上升情形给出 \(\ell(szt)=n+1\)，而 \(szt=wt\) 仅长 \(n-1\)，矛盾。因此 \(\ell(swt)=n-2\)。将这四种长度关系代入算子定义，两种复合的共同值为
\[
\begin{array}{c|c|l}
(\ell(sw)-n,\ell(wt)-n)&\ell(swt)&L_sR_te_w=R_tL_se_w\\\hline
(1,1)&n+2&e_{swt}\\
(1,-1)&n&e_{swt}+de_{sw}\\
(-1,1)&n&e_{swt}+de_{wt}\\
(-1,-1)&n-2&e_{swt}+d(e_{sw}+e_{wt})+d^2e_w
\end{array}
\]
若 \(sw=wt\)，两次长度改变方向相同。都增加时，两种复合均为 \(e_w+de_{sw}\)；都减少时，均为 \((1+d^2)e_w+de_{sw}\)。于是左右算子确实交换。

当 \(m_{st}<\infty\) 时，\(s,t\) 的两条长度 \(m_{st}\) 交替词是二面体子群 \(W_{\{s,t\}}\) 最长元的约化词。由\cref{b5:thm:coxeter-word-properties}证明中所用的标准抛物子群长度相容性，子群内的长度等于全群中的长度，故这两个词在 \(W\) 中也约化，这里调用的是该证明中 \(W_Iw\) 最短代表分解前明确列出的长度结论。两个交替词的每个后缀都约化，左算子从右向左作用时每步都增加长度，故不产生 \(d\) 项；在 \(e_1\) 上的共同值就是该二面体最长元的 \(e_{w_0}\)。两种乘积都与全部 \(R_u\) 交换，而对约化词 \(w=s_1\cdots s_r\)，依次右乘给出
\[
e_w=R_{s_r}\cdots R_{s_1}e_1.
\]
右边最先作用的是 \(R_{s_1}\)，每个词的前段都约化，故每一步都不产生 \(d\) 项。于是两种左算子乘积在全 \(E\) 上相同。这样 \(L_s\) 满足全部定义关系，给出 \(\mathcal H(W)\) 在 \(E\) 上的表示。最后 \(H_we_1=e_w\)，所以任何 \(H_w\) 的线性关系作用到 \(e_1\) 后都只能有零系数。
\end{proof}

原始 Hecke 代数的基及乘法约定见\cite[Section 1, p. 165]{KazhdanLusztig1979Hecke}。原文使用 \(q=v^2\) 与 \(T_w=v^{\ell(w)}H_w\)，于是二次式变为 \((T_s+1)(T_s-q)=0\)。比较公式时须先写清参数与归一化的对应。

\ProseParagraphBreak
在秩一情形无需任何一般词理论：
\[
\mathcal H(C_2)=R[H]/\bigl(H^2-(v-v^{-1})H-1\bigr).
\]
首一除法直接给出基 \(1,H\)。当 \(v=1\) 时，关系成为 \(H^2=1\)，恢复群环 \(\mathbb Z[C_2]\)；一般 Coxeter 情形由标准基定理同样得到
\(\mathcal H(W)/(v-1)\simeq\mathbb Z[W]\)。
\ProseParagraphBreak
在 \(S_3=\langle s,t\rangle\) 中，标准基为
\[
1,\ H_s,\ H_t,\ H_{st},\ H_{ts},\ H_{sts},
\]
而 \eqref{b5:eq:hecke-standard-multiplication}给出
\[
H_sH_{st}=H_t+(v-v^{-1})H_{st},\qquad
H_sH_{ts}=H_{sts}.
\]
一个乘法用了长度下降，另一个用了长度上升；不能在两种情况下都直接把群元素下标相乘。

\subsection{反演算子与规范基}
\label{b5:subsec:A04:02}

二次关系给出 \(H_s^{-1}=H_s-(v-v^{-1})\)。由此可定义半线性环对合
\[
\overline v=v^{-1},\qquad
\overline{H_w}=H_{w^{-1}}^{-1}.
\]
它是 \(v\mapsto v^{-1}\)、\(H_s\mapsto H_s^{-1}\) 的乘法延拓。沿约化词 \(w=s_1\cdots s_r\)，有 \(\overline{H_w}=H_{s_1}^{-1}\cdots H_{s_r}^{-1}=(H_{s_r}\cdots H_{s_1})^{-1}=H_{w^{-1}}^{-1}\)，故保持乘法而反演系数参数。良定义可从生成关系核验：逆元满足参数反演后的二次关系，取逆并反向读取辫关系仍得到同一对辫词；再作一次变换恢复各生成元。例如 \(\overline{H_{st}}=H_s^{-1}H_t^{-1}\)，而 \(H_{st}^{-1}=H_t^{-1}H_s^{-1}\)，一般不同。因此 bar 不能与取逆所造成的乘法次序倒转混用。

\begin{theorem}[Kazhdan--Lusztig 规范基]\label{b5:thm:kl-basis}
设 \((W,S)\) 为 Coxeter 系统，\(\ell\) 为关于 \(S\) 的长度函数，\(v\) 为形式变量，\(R=\mathbb Z[v,v^{-1}]\)。令 \(\mathcal H(W)\) 为 \(R\) 上的等参数 Hecke 代数，满足 \((H_s-v)(H_s+v^{-1})=0\) 及 Coxeter 辫关系，\(H_w\) 为其标准基。以 \(\overline v=v^{-1}\)、\(\overline{H_w}=H_{w^{-1}}^{-1}\) 定义半线性环对合。对每个 \(w\in W\)，存在唯一元素 \(C'_w\in\mathcal H(W)\) 满足
\[
\overline{C'_w}=C'_w,\qquad
C'_w=H_w+\sum_{y<w}a_{y,w}(v)H_y,\qquad
a_{y,w}(v)\in v^{-1}\mathbb Z[v^{-1}],
\]
其中 \(<\) 为 Bruhat 严格序。这些元素构成 \(\mathcal H(W)\) 的 \(R\) 基，且对形式变量 \(q\) 有
\[
a_{y,w}(v)=v^{\ell(y)-\ell(w)}\cdot P_{y,w}(v^2),
\qquad P_{y,w}(q)\in\mathbb Z[q],
\]
对 \(y<w\) 有 \(\deg P_{y,w}\le\bigl(\ell(w)-\ell(y)-1\bigr)/2\)，并置 \(P_{w,w}=1\)。
\end{theorem}
\symidx{\(C'_w,P_{y,w}(q)\)：Kazhdan--Lusztig 规范基元与多项式}

\begin{proof}
先注意 \(\overline{H_w}=H_w+\sum_{y<w}r_{y,w}H_y\)：沿约化词把每个 \(H_s\) 换成 \(H_s-d\)，展开 \(\prod_i(H_{s_i}-d)\)，每项对应保留下来的位置子集。使用乘法规则时，要么接上该字母，要么在长度下降时留在原下标；最后得到的元素仍是原词某个子词的乘积，故由子词判据属于 \([1,w]\)。真位置子集少于 \(\ell(w)\) 个字母，不能产生长度为 \(\ell(w)\) 的 \(w\)；因此最高项只来自全部位置，其系数为一。由\cref{b5:thm:coxeter-word-properties}的第三项，主下区间 \([1,w]\) 有限，因此下面对固定 \(w\) 的各个和式均为有限和。

按长度构造 \(C'_w\)，取 \(C'_1=1\)。设 \(w=su\) 且 \(\ell(w)=\ell(u)+1\)，较短元素的规范基已经构造。乘积
\(D=(H_s+v^{-1})C'_u\) 自对偶。最高项 \(H_u\) 的贡献为 \(H_w+v^{-1}H_u\)；对 \(y<u\)，乘法规则给出
\[
(H_s+v^{-1})a_{y,u}H_y=
\begin{cases}
a_{y,u}H_{sy}+v^{-1}a_{y,u}H_y,&sy>y,\\
a_{y,u}H_{sy}+v a_{y,u}H_y,&sy<y.
\end{cases}
\]
子词判据保证这里的下标 \(y,sy\) 都小于 \(w\)。各个 \(a_{y,u}\) 只含严格负次幂，故上式的系数都没有正次幂；常数项只可能来自 \(sy<y\) 时的 \(v a_{y,u}\)，并等于 \(a_{y,u}\) 中 \(v^{-1}\) 的系数 \(\mu(y,u)\)。因此定义
\[
 C'_w=D-\sum_{\substack{y<u\\sy<y}}\mu(y,u)C'_y.
\]
减项仍自对偶，恰消去各个常数项；其低项系数原本只有负次幂，故不会产生新的常数项。这给出所需存在性。

同时归纳验证：\(a_{y,w}\) 中每个指数不小于 \(\ell(y)-\ell(w)\)，且与该数同奇偶。设 \(a_{y,u}\) 的一个指数为 \(k\)，则归纳假设给出
\[
k\ge\ell(y)-\ell(u),\qquad
k\equiv\ell(y)-\ell(u)\pmod2.
\]
在 \(sy>y\) 的情形中，\(H_{sy}\) 的指数 \(k\) 及 \(H_y\) 的指数 \(k-1\) 分别不小于 \(\ell(sy)-\ell(w)\) 及 \(\ell(y)-\ell(w)\)，奇偶也分别一致。在 \(sy<y\) 的情形中，\(H_{sy}\) 的指数 \(k\) 及 \(H_y\) 的指数 \(k+1\) 的下界分别比上述目标下界多二，故同样满足要求。最高项的额外贡献 \(v^{-1}H_u\) 也满足这一约束。

再检查减项。为统一记号，置 \(a_{y,y}=1\)。若 \(\mu(y,u)\ne0\)，则 \(-1\equiv\ell(y)-\ell(u)\pmod2\)，所以 \(\ell(w)-\ell(y)\) 为偶数。对每个 \(x\le y\)，\(C'_y\) 中 \(H_x\) 的系数的指数下界为
\[
\ell(x)-\ell(y)\ge\ell(x)-\ell(w),
\qquad
\ell(x)-\ell(y)\equiv\ell(x)-\ell(w)\pmod2.
\]
这也包括 \(x=y\) 时的指数零。因此减项保持所需的下界与奇偶性。结合 \(C'_w\) 的所有下位系数只含严格负次幂，便可唯一写成
\(v^{\ell(y)-\ell(w)}P_{y,w}(v^2)\)。若 \(P_{y,w}\) 中出现 \(q^j\)，则 \(\ell(y)-\ell(w)+2j\le-1\)，这正给出所述次数上界。

若有两个符合刻画的元素，取其差中 Bruhat 极大的非零下标 \(y\)。三角形反演说明对应系数 \(f\) 必满足 \(f(v)=f(v^{-1})\)。但 \(f\) 只有严格负次幂，只能为零，矛盾。这证明唯一性。最后，对每个 \(w\)，在有限主下区间 \([1,w]\) 上反解单位三角形变换，即可把 \(H_w\) 写成有限个 \(C'_y\) 的线性组合；同一三角形关系也给出线性无关性，故全部 \(C'_w\) 组成 \(R\) 基。
\end{proof}

这是原论文\cite[Theorem 1.1, equivalent formulation (1.1.c), p. 166]{KazhdanLusztig1979Hecke}的归一化改写。多项式 \(P_{y,w}\) 称为\term[Kazhdan-Lusztig-duoxiangshi]{Kazhdan--Lusztig 多项式}[Kazhdan--Lusztig polynomial]。下述低秩计算可直接检验递推；一般系数的正性以及 Lie 代数中的特征标公式则是另外的结果。

\ProseParagraphBreak
对简单反射，
\[
C'_s=H_s+v^{-1}.
\]
由逆元公式，\(\overline{H_s+v^{-1}}=H_s-(v-v^{-1})+v=H_s+v^{-1}\)，所以确为自对偶。若先只要求 \(H_s+a(v)\) 自对偶，条件其实是
\[
a(v)-a(v^{-1})=v^{-1}-v.
\]
除了解 \(a(v)=v^{-1}\)，还可加上任意满足 \(b(v)=b(v^{-1})\) 的 Laurent 多项式。要求 \(a(v)\) 只含严格负次幂，便排除了这份自由度：两解的差同时在参数反演下不变且只含负次幂，只能为零。这正是次数条件的规范化作用。若记 \(a=v^{-1}\)，在 \(S_3\) 中得到
\begin{align*}
C'_{st}&=C'_sC'_t
=H_{st}+a(H_s+H_t)+a^2,\\
C'_{ts}&=C'_tC'_s
=H_{ts}+a(H_s+H_t)+a^2,\\
C'_{sts}&=C'_sC'_tC'_s-C'_s\\
&=H_{sts}+a(H_{st}+H_{ts})+a^2(H_s+H_t)+a^3.
\end{align*}
每个乘积和差均自对偶；展开后最高项系数一，其余系数只含严格负次幂，且下标在已计算的 Bruhat 区间中。因此它们就是所求规范基。最后一个式子的减项不可省略：未减之前 \(H_s\) 的系数为 \(1+a^2\)，含常数项，不满足规范条件。这里所有 \(P_{y,w}\)（\(y\le w\)）都等于一，但高秩时不会一直如此。

\ProseParagraphBreak
例如取前面已经算过的 \(w=3412=s_2s_1s_3s_2\in S_4\)，其中 \(s_i=(i,i+1)\)、\(\ell(w)=4\)。仍记 \(a=v^{-1}\)，令 \(u=s_1s_3s_2\)。它的三个字母互不相同，逐次相乘没有产生常数修正项，得到
\begin{align*}
C'_u={}&C'_{s_1}C'_{s_3}C'_{s_2}\\
={}&H_u+a(H_{s_1s_3}+H_{s_1s_2}+H_{s_3s_2})\\
&+a^2(H_{s_1}+H_{s_3}+H_{s_2})+a^3.
\end{align*}
在这里的八个下标中，左乘 \(s_2\) 使长度下降的只有 \(s_2\) 本身。它在 \(C'_u\) 中的系数为 \(a^2\)，没有 \(v^{-1}\) 项，因此规范基递推的减项为零。再用 \(H_{s_2}^2=1+(v-v^{-1})H_{s_2}\) 展开 \(C'_{s_2}C'_u\)：下标 \(s_2\) 的系数为
\[
a^2(v-v^{-1})+2a^3=a+a^3,
\]
单位元的系数为 \(a^2+a^4\)。求和遍历整个 Bruhat 区间 \([1,w]\)；子词判据给出其中长度 \(1,2,3\) 的元素数分别为 \(3,5,4\)。除 \(1,s_2\) 外，各项系数仍为 \(a^{4-\ell(y)}\)。和式已给 \(s_2\) 系数 \(a^3\) 及单位元系数 \(a^4\)，故另外添加 \(aH_{s_2}+a^2\) 恰补出刚算得的两项；合起来是
\[
C'_w=C'_{s_2}C'_u
=\sum_{y\le w}a^{4-\ell(y)}H_y+aH_{s_2}+a^2.
\]
所有下位系数都只有严格负次幂，所以这正是 \(C'_w\)。特别地，
\[
a_{1,w}(v)=v^{-4}+v^{-2},\qquad P_{1,w}(q)=1+q.
\]
因此 \(\deg P_{1,w}=1\)，达到整数次数所能取的上界 \(\lfloor(4-1)/2\rfloor=1\)。多出来的 \(q\) 项不能仅从下标的长度归一化中读出。原论文也列出了这个 \(A_3\) 型算例，见\cite[Section 6, Example 6.1, p. 177]{KazhdanLusztig1979Hecke}。

\ProseParagraphBreak
规范基的系数属于 \(\mathbb Z[v,v^{-1}]\)，变换矩阵是单位三角的。因此取任意域 \(k\) 和 \(z\in k^\times\)，令 \(v=z\) 后，标准基和规范基仍分别是所得 \(k\) 代数的基；反解基变换不用除以可能在 \(z\) 处消失的多项式。这种稳定的坐标并不保证特殊化后的代数半单，秩一已经能看出区别。

\begin{proposition}\label{b5:prop:hecke-rank-one-specialization}
设 \(R=\mathbb Z[v,v^{-1}]\)，
\[
\mathcal H(C_2)=R[H]/\bigl(H^2-(v-v^{-1})H-1\bigr),
\qquad C'=H+v^{-1}.
\]
则 \(1,C'\) 为 \(R\) 基，且 \((C')^2=(v+v^{-1})C'\)。在 \(\mathbb Q(v)\) 上，
\[
e_+=\frac{C'}{v+v^{-1}},\qquad
e_-=\frac{v-H}{v+v^{-1}}
\]
为两个正交本原幂等元，\(e_++e_-=1\)。在 \(v=i\) 处特殊化到 \(\mathbb C\) 后，\(1,C'\) 仍为基，但 \(C'=H-i\ne0\)、\((C')^2=0\)，所得代数同构于 \(\mathbb C[\epsilon]/(\epsilon^2)\)，是非半单局部代数；上述两个本原幂等元的分式在该参数处不能定义。
\end{proposition}
\begin{proof}
首一二次式的除法给出基 \(1,H\)，而 \(H=C'-v^{-1}\)，所以 \(1,C'\) 也是基。二次关系给出
\[
(H+v^{-1})^2
=1+(v+v^{-1})H+v^{-2}
=(v+v^{-1})(H+v^{-1}).
\]
在 \(\mathbb Q(v)\) 中，两个根 \(v\)、\(-v^{-1}\) 不同，求值到这两个根给出
\[
\mathcal H(C_2)\otimes_R\mathbb Q(v)
\simeq\mathbb Q(v)\times\mathbb Q(v).
\]
在这个同构下，\(e_+\)、\(e_-\) 分别对应 \((1,0)\)、\((0,1)\)，故具有所述幂等性质且本原。

令 \(v=i\)，二次关系成为 \((H-i)^2=0\)。\(1,H-i\) 仍线性无关，因此 \(\epsilon=H-i\) 是一个非零的平方零元，并给出所述双数代数。元素 \(a+b\epsilon\) 可逆当且仅当 \(a\ne0\)：此时逆为 \(a^{-1}-a^{-2}b\epsilon\)，而 \(a=0\) 时它幂零。因此唯一极大理想是非零的 \((\epsilon)\)，该代数是局部而非半单的。两个幂等元分式实际已定义在 \(R[(v+v^{-1})^{-1}]\otimes_R\mathcal H(C_2)\) 中。\(v\mapsto i\) 把被求逆的元素送到零，故 \(R\to\mathbb C\) 不能延拓到这个局部化环，分式无法沿该映射特殊化；原代数无需这些分母，仍可特殊化。
\end{proof}

在 \(v=1\) 处特殊化到 \(\mathbb Q\)，分母变成二，可以求逆；两个幂等元成为 \((1+H)/2\)、\((1-H)/2\)，分别投影到 \(\mathbb Q[C_2]\) 的平凡和符号部分。

这与\chapref{b5:ch:08}中的模特征现象相似：改变参数或特征，会使原先可用幂等元分开的成分合并，并出现非平凡扩张。规范基提供代数元素的坐标，本原幂等元用来分解正则模；二者承担的任务不同。

\subsection{胞腔与一个二维表示}
\label{b5:subsec:A04:03}

\begin{definition}\label{b5:def:kl-left-cell}
设 \((W,S)\) 为 Coxeter 系统，\(v\) 为形式变量，\(\mathcal H(W)\) 为 \(\mathbb Z[v,v^{-1}]\) 上的等参数 Hecke 代数，满足 \((H_s-v)(H_s+v^{-1})=0\)，\(C'_w\)、\(w\in W\)，为 Kazhdan--Lusztig 规范基。对 \(x,y\in W\)，若某个 \(s\in S\) 使 \(C'_sC'_y\) 的展开含非零的 \(C'_x\) 系数，就画有向关系 \(y\to x\)。令 \(x\preceq_L y\) 表示存在从 \(y\) 到 \(x\) 的有向链，允许空链。互相可达给出的等价类称为\term[zuo-baoqiang]{左胞腔}[left cell]。改用右乘定义右胞腔，同时允许两种乘法则定义双边胞腔。
\end{definition}
\symidx{\(x\preceq_L y\)：Hecke 规范基乘法定义的左胞腔预序}

这不是 Bruhat 偏序：它由乘法结构常数决定，互相可达还可能把不同长度的元素放在一起。定义采用的规范基版本必须一致；原文同时使用 \(C_w,C'_w\)，其等价的图定义及表示论用途见\cite[Section 1, Definition 1.2 and Theorems 1.3--1.4, pp. 166--167]{KazhdanLusztig1979Hecke}。

\ProseParagraphBreak
令 \(\Gamma\) 是一个左胞腔。所有满足 \(x\preceq_L y\)、某个 \(y\in\Gamma\) 的基元的线性生成空间是左理想，因为乘以 \(C'_s\) 只能沿有向关系继续向下走，而 \(H_s=C'_s-v^{-1}\)。其中不属于 \(\Gamma\) 的那些基元也生成左理想：若还能乘回 \(\Gamma\)，它就与 \(\Gamma\) 互相可达，矛盾。这两个左理想的商因而给出一个以 \(\Gamma\) 为基的胞腔表示。这一步只用定义与基的线性无关性。

\ProseParagraphBreak
在 \(S_3\) 中，置 \(b=v+v^{-1}\)，直接由前面的展开和二次关系得到
\[
(C'_s)^2=bC'_s,\quad
C'_tC'_s=C'_{ts},\quad
C'_sC'_{ts}=C'_{sts}+C'_s,\quad
C'_tC'_{ts}=bC'_{ts}.
\]
交换 \(s,t\) 给出另外四式，其中 \(tst=sts\) 仍是同一个最长元；而 \(C'_sC'_{sts}=C'_tC'_{sts}=bC'_{sts}\)。忽略自环，按乘积中出现的下标读出箭头 \(1\to s,t\)，\(s\to ts\)、\(ts\to s,sts\)，\(t\to st\)、\(st\to t,sts\)。只有 \(s,ts\) 与 \(t,st\) 分别能互相到达，其余两点单独成类；由此左胞腔恰为
\[
\{1\},\qquad\{s,ts\},\qquad\{t,st\},\qquad\{sts\}.
\]
对 \(\Gamma=\{s,ts\}\)，上述两个左理想具体为
\[
I=R C'_s\oplus R C'_{ts}\oplus R C'_{sts},\qquad
J=R C'_{sts}.
\]
仅由 \(C'_s,C'_{ts}\) 张成的空间并不稳定，因为 \(C'_sC'_{ts}\) 含有 \(C'_{sts}\)。所取的是商模 \(I/J\)，令其中两基向量为 \(e_s=C'_s+J,e_{ts}=C'_{ts}+J\)。因此 \(C'_{sts}\) 项在商中为零，得到
\[
[C'_s]_\Gamma=
\begin{pmatrix}b&1\\0&0\end{pmatrix},
\qquad
[C'_t]_\Gamma=
\begin{pmatrix}0&0\\1&b\end{pmatrix}.
\]
由 \(H_i=C'_i-v^{-1}\) 得到 Hecke 生成元矩阵。在 \(v=1\) 处，它们成为
\[
s\longmapsto\begin{pmatrix}1&1\\0&-1\end{pmatrix},
\qquad
t\longmapsto\begin{pmatrix}-1&0\\1&1\end{pmatrix}.
\]
两者平方为一，乘积迹为 \(-1\)、行列式为一，且满足辫关系。三个换位的迹为零，两个三轮换的迹为 \(-1\)，所以该胞腔表示复化后与 \(S_3\) 的二维标准表示同构。胞腔方法由此把一个乘法基中的组合结构变成了实际的表示矩阵。

\ProseParagraphBreak
这些胞腔也能在前面的插入算法中认出。按通常的置换复合与单行记号，\(s=213\)、\(ts=312\)，\(t=132\)、\(st=231\)。对四个词分别行插入，新增格的位置给出
\[
Q(s)=Q(ts)=\begin{array}{cc}1&3\\2&\end{array},
\qquad
Q(t)=Q(st)=\begin{array}{cc}1&2\\3&\end{array}.
\]
另外，\(1=123\) 的记录表是一行 \(1,2,3\)，\(sts=321\) 的记录表是一列 \(1,2,3\)。因此在 \(S_3\) 中，左胞腔恰好是固定记录表 \(Q\) 的置换集合；改用右胞腔则固定插入表 \(P\)，双边胞腔只保留共同的形状。对一般 \(S_n\)，这些对应仍成立，其证明还要把胞腔等价与 Knuth 的表等价联系起来，见\cite[Section 5, pp. 176--177]{KazhdanLusztig1979Hecke}。这里的插入计算只验证了 \(S_3\) 的情形。

\ProseParagraphBreak
从记录表的形状识别胞腔模的 Specht 类型时，还须保留本节使用 \(C'_w\) 的约定。对胞腔 \(\{1\}\)，\(C'_s,C'_t\) 的像都为零，所以 \(H_s,H_t\) 都按 \(-v^{-1}\) 作用；在 \(v=1\) 时得到符号表示。对 \(\{sts\}\)，两者都按 \(b\) 作用，\(H_s,H_t\) 则按 \(v\) 作用，特殊化后得到平凡表示。相应的 RS 形状却分别为 \((3)\)、\((1,1,1)\)，而这两种形状的 Specht 模分别是平凡表示与符号表示。因此不能在没有转换规范基约定的情况下，把 RS 形状 \(\lambda\) 直接写成胞腔模的 Specht 标号 \(\lambda\)。两个二维胞腔的形状 \((2,1)\) 与其共轭分拆相同，单看这个例子不会显出差别。


\begin{chaptersummary}
邻位换位每次改变一个逆序，所以最短词与逆序数一致。Bruhat 序允许乘任意反射，弱序只允许简单生成元；两者从 \(S_3\) 开始已有区别。半标准表通过逐次加横条解释 Schur 多项式，而行插入与逆插入把置换分配到同形状的成对标准表，给出标准表数平方和的恒等式。

\ProseParagraphBreak
轨道计数利用轨道大小与稳定子群阶的关系，子空间计数则把线性无关元组的数目除以每个子空间的有序基数。生成函数记录次数、逆序数或颜色数量。Hecke 代数将反射的二次关系变形；它的规范基以自对偶和三角条件刻画，低秩胞腔已经能产生非平凡表示矩阵。
\end{chaptersummary}

\BookExercises

\begin{exercise}\label{b5:ex:A-dihedral-four}
设 \(W=\bigl\langle s,t\mid s^2=t^2=1,(st)^4=1\bigr\rangle\)。列出全部元素的长度，并求长度生成多项式。
\end{exercise}
\begin{solution}
元素依次为
\[
1;\quad s,t;\quad st,ts;\quad sts,tst;\quad stst=tsts.
\]
在长度小于四时，约去相邻相同字母后只剩交替词，两个不同交替词在正方形的旋转反射模型中不同。由\subsecref{b5:subsec:A01:01}的秩二模型，所列八个元素确实互异。长度四的等式 \(stst=tsts\) 来自 \((st)^4=1\)，而非长度三的辫关系；其他更长词可利用这个关系并约去相同邻字继续缩短。因而所列长度分别为 \(0,1,2,3,4\)，生成多项式为 \(1+2q+2q^2+2q^3+q^4\)。
\end{solution}

\begin{exercise}\label{b5:ex:A-bubble-word}
求置换 \(4213\) 的长度，并用邻位换位写出一个约化表达式。
\end{exercise}
\begin{solution}
四贡献三个逆序，二与一再贡献一个，所以长度四。从 \(4213\) 右乘 \(s_2,s_1,s_2,s_3\)，依次得到 \(4123,1423,1243,1234\)。因此
\(4213=s_3s_2s_1s_2\)。这个词有四项，达到逆序下界，故约化。
\end{solution}

\begin{exercise}\label{b5:ex:A-bruhat-not-length}
在 \(S_3\) 中给出两个长度相同却不能在 Bruhat 序下比较的不同元素；再给出一对在 Bruhat 序下可比而在右弱序下不可比的元素，并直接核验。
\end{exercise}
\begin{solution}
取 \(s=(12)\)、\(t=(23)\)，它们均长度一。Bruhat 序中 \(u<v\) 必有 \(\ell(u)<\ell(v)\)，所以这两个不同元素不可比。另一例取 \(t\) 与 \(st\)：右乘反射 \(tst\) 有 \(t(tst)=st\)，长度增加，故 \(t\le st\)。右弱序也严格增加长度，故唯一可能的方向是 \(t\le_R st\)。长度差一迫使 \(st=tr\)、\(r\in\{s,t\}\)；但 \(r=s\) 给出错误的 \(st=ts\)，\(r=t\) 则给出 \(st=1\)，都不成立。
\end{solution}

\begin{exercise}\label{b5:ex:A-longest-complement}
设 \(w_0=n(n-1)\cdots1\in S_n\)。证明
\(\ell(w_0w)=\binom n2-\ell(w)\)，并据此说明长度多项式的系数为何对称。
\end{exercise}
\begin{solution}
\((w_0w)(i)=n+1-w(i)\)，所以每一对位置在 \(w_0w\) 中逆序，当且仅当在 \(w\) 中不逆序。位置对共有 \(\binom n2\)，由长度等于逆序数得到公式。映射 \(w\mapsto w_0w\) 是双射，因此长度 \(d\) 与长度 \(\binom n2-d\) 的置换等数。
\end{solution}

\begin{exercise}\label{b5:ex:A-semistandard-21}
求用 \(1,2,3\) 填写的形状 \((2,2)\) 的半标准表数，按重量写出 \(s_{(2,2)}(x_1,x_2,x_3)\)，并说明对应的权空间维数。
\end{exercise}
\begin{solution}
每列自上而下只能是 \((1,2),(1,3),(2,3)\) 之一。它们按上下两个分量分别比较，依次递增；行弱递增要求左列不大于右列。因此两列的选择为前三种列中可重复的无序二元组，共六种。两列相同时，重量依次为 \((2,2,0),(2,0,2),(0,2,2)\)；两列不同时，重量为 \((2,1,1),(1,2,1),(1,1,2)\)。每种重量恰有一张表，故
\[
s_{(2,2)}(x_1,x_2,x_3)
=\sum_{i<j}x_i^2x_j^2+x_1x_2x_3(x_1+x_2+x_3).
\]
在特征零的 Schur 模中，这六个重量对应的权空间各为一维，总维数为六。这里没有不同表给出相同重量的情形，可与\cref{b5:prop:A-schur-21-weights}中形状 \((2,1)\) 的二重权空间比较。
\end{solution}

\begin{exercise}\label{b5:ex:A-insert-2413}
对置换 \(35124\) 进行行插入，求 \(P,Q\)，再从所得表对逆向恢复置换。计算逆排列的表对，并核验形状与最长递增、递减子列长度的关系。
\end{exercise}
\begin{solution}
前两步第一行为 \(3,5\)。插入一时弹出三，得到第一行 \(1,5\)、第二行 \(3\)；插入二时弹出五，得到第一行 \(1,2\)、第二行 \(3,5\)；最后把四附到第一行。因此
\[
P=\begin{array}{ccc}1&2&4\\3&5&\end{array},
\qquad Q=\begin{array}{ccc}1&2&5\\3&4&\end{array}.
\]
逆插入先删记录数字五所在的第一行末格，输出四；再从第二行取出五，在第一行替换二，输出二；随后从第二行取出三，替换第一行的一，输出一。余下第一行 \(3,5\)，依次输出五、三。逆序读取便恢复 \(3,5,1,2,4\)。

逆排列为 \(34152\)。直接插入得到
\[
P'=\begin{array}{ccc}1&2&5\\3&4&\end{array},
\qquad Q'=\begin{array}{ccc}1&2&4\\3&5&\end{array},
\]
确实交换了原表对。子列 \(1,2,4\) 递增；含开头的 \(3\) 时，后面只有 \(5,4\) 比它大，且不能依此顺序同时取二者；含 \(5\) 时更无法接上后面的数，故这样的递增子列至多长二。不含 \(3,5\) 时，只剩递增的 \(1,2,4\)，故最长长三。递减子列 \(5,1\) 长二；开头 \(3,5\) 与末尾 \(1,2,4\) 分别递增，递减子列至多从每组取一个，故长度不能超过二。因此最长递减子列长二，与形状 \((3,2)\) 的第一行和第一列长度一致。
\end{solution}

\begin{exercise}\label{b5:ex:A-inverse-pair}
给定
\[
P=\begin{array}{ccc}1&3&4\\2&&\end{array},
\qquad
Q=\begin{array}{ccc}1&2&4\\3&&\end{array},
\]
求对应的置换。
\end{exercise}
\begin{solution}
删记录数字四的格子直接输出四。删数字三所在的第二行格子，取出二，在上行把一替换为二，输出一。此时上行为 \(2,3\)，接着输出三、二。因此置换为 \(2314\)。正向插入二、三后第一行为 \(2,3\)，插入一弹出二，最后把四附到第一行，核验了原表对。
\end{solution}

\begin{exercise}\label{b5:ex:A-catalan-five}
列出形状 \((3,3)\) 的全部标准表，只写第一行即可，并与 Catalan 数核验。
\end{exercise}
\begin{solution}
第一行可能是 \(123,124,125,134,135\)，第二行由其余三个数字递增填入。按 \(1,2,\ldots,6\) 的顺序读取数字，对每个前缀 \(\{1,\ldots,j\}\)，它们都满足其中第一行的数字数不少于第二行的数字数；这等价于第一行第 \(i\) 个数字小于第二行第 \(i\) 个数字，见\subsecref{b5:subsec:A02:03}中的计数论证。其他选择会违反该条件。共有五张，符合 \(\dfrac14\binom63=5\)。
\end{solution}

\begin{exercise}\label{b5:ex:A-necklace-three-colors}
四个圆周位置用三种固定颜色染色，只把旋转视为相同。求总轨道数，并求其中第一种颜色出现两次、其余颜色各一次的轨道数。
\end{exercise}
\begin{solution}
循环指标给出总数
\((3^4+3^2+2\cdot3)/4=24\)。
指定重量为 \(a^2bc\) 时，只有恒等置换的固定染色有贡献，系数为
\[
\frac14\cdot\frac{4!}{2!1!1!}=3.
\]
半周旋转强迫每种颜色出现偶数次，四分之一周旋转强迫全部同色，故其余项不贡献。
\end{solution}

\begin{exercise}\label{b5:ex:A-bracelet-three-colors}
把上一题的等价关系扩大为正方形的全部旋转和反射，求三色染色的轨道数。
\end{exercise}
\begin{solution}
四个旋转的固定染色总和为 \(96\)。两条经过对顶点的反射轴具有循环型 \(1^22\)，各固定 \(3^3=27\) 种；两条经过对边中点的轴具有型 \(2^2\)，各固定 \(3^2=9\) 种。因此总轨道数为
\((96+54+18)/8=21\)。这小于只允许旋转时的二十四，说明多三对旋转轨道被反射合并。
\end{solution}

\begin{exercise}\label{b5:ex:A-partition-six}
分别列出六的不同部分分拆与全部部分为奇数的分拆，并用正文的二进制双射将它们逐一配对。
\end{exercise}
\begin{solution}
不同部分分拆为 \(6,5+1,4+2,3+2+1\)。把每个部分写成二的幂乘奇数，便依次得到
\[
\begin{aligned}
6&\longleftrightarrow3+3,&
5+1&\longleftrightarrow5+1,\\
4+2&\longleftrightarrow1+1+1+1+1+1,&
3+2+1&\longleftrightarrow3+1+1+1.
\end{aligned}
\]
右边正是六的全部奇数部分分拆。反向将每个奇数的出现次数写成二进制，就恢复左边的不同部分分拆。
\end{solution}

\begin{exercise}\label{b5:ex:A-length-four-coefficient}
求 \(S_4\) 中长度为二的置换数，不必逐个列出置换。
\end{exercise}
\begin{solution}
长度多项式为 \((1+q)(1+q+q^2)(1+q+q^2+q^3)\)。按三个因子的顺序，取得 \(q^2\) 时所选项的次数分别为
\[
(0,0,2),\quad(0,2,0),\quad(1,1,0),\quad(1,0,1),\quad(0,1,1).
\]
每种选择贡献一，故 \(q^2\) 的系数为五。
\end{solution}

\begin{exercise}\label{b5:ex:A-gaussian-four-two}
在 \(\mathbb F_2^4\) 中固定两条不同直线 \(L_1,L_2\)。求包含 \(L_1\) 而不包含 \(L_2\) 的二维子空间数，以及两条直线都不包含的二维子空间数。
\end{exercise}
\begin{solution}
\cref{b5:prop:gaussian-subspace-count}给出的二维子空间总数为 \({4\brack2}_2=35\)。包含 \(L_i\) 的二维空间与 \(\mathbb F_2^4/L_i\) 的直线一一对应，故各有 \((2^3-1)/(2-1)=7\) 个。商空间中每条直线的原像维数为二且包含 \(L_i\)；反过来，每个包含 \(L_i\) 的二维空间取商后是一条直线。取原像与取商互为逆，故确为一一对应。同时包含 \(L_1,L_2\) 的二维空间只有 \(L_1+L_2\) 一个。因此第一问为 \(7-1=6\)；第二问由容斥为 \(35-7-7+1=22\)。
\end{solution}

\begin{exercise}\label{b5:ex:A-flag-three}
计算 \(\mathbb F_3^3\) 的完整旗数，并用“先选直线，再选平面”重新计算。
\end{exercise}
\begin{solution}
公式给出 \((1)(1+3)(1+3+9)=52\)。直接计算时，直线共有十三条；固定直线 \(L\) 后，包含它的平面对应二维商空间的直线，共四条。因此也得到 \(13\cdot4=52\)。
\end{solution}

\begin{exercise}\label{b5:ex:A-hecke-rank-one-powers}
在秩一 Hecke 代数中，令 \(d=v-v^{-1}\)。计算 \(H^3,H^4\)，并核验 \(v=1\) 时的结果。
\end{exercise}
\begin{solution}
由 \(H^2=dH+1\)，
\[
H^3=(d^2+1)H+d,\qquad
H^4=(d^3+2d)H+d^2+1.
\]
令 \(v=1\) 后 \(d=0\)，分别得到 \(H^3=H,H^4=1\)，符合二阶群生成元的幂。
\end{solution}

\begin{exercise}\label{b5:ex:A-hecke-inverse}
在 \(S_3\) 的 Hecke 代数中计算 \(H_{st}^{-1}\)，并解释为什么它不同于 \(\overline{H_{st}}\)。
\end{exercise}
\begin{solution}
令 \(d=v-v^{-1}\)。则
\[
H_{st}^{-1}=H_t^{-1}H_s^{-1}
=H_{ts}-d(H_s+H_t)+d^2.
\]
而 \(\overline{H_{st}}=H_s^{-1}H_t^{-1}
=H_{st}-d(H_s+H_t)+d^2\)。标准基中 \(H_{st}\ne H_{ts}\)，所以两者不同。反演参数的横线映射保持乘法次序，取群元素逆的运算则反转次序。
\end{solution}

\begin{exercise}\label{b5:ex:A-kl-simple-square}
设 \(k\) 为域，将秩一 Hecke 代数沿 \(v\mapsto1\) 特化为 \(A_k=k[H]/(H^2-1)\)。比较 \(\operatorname{char}k\ne2\) 与 \(\operatorname{char}k=2\) 时的代数分解，分别求本原幂等元和 Jacobson 根，并解释参数同为一为何仍可能得到不同的半单性结论。
\end{exercise}
\begin{solution}
若 \(\operatorname{char}k\ne2\)，两个根 \(1,-1\) 不同，中国剩余定理给出 \(A_k\simeq k\times k\)。两个本原幂等元为
\[
e_+=\frac{1+H}{2},\qquad e_-=\frac{1-H}{2},
\]
Jacobson 根为零。

若 \(\operatorname{char}k=2\)，令 \(\varepsilon=H+1\)，则 \(\varepsilon^2=0\)，而 \(1,\varepsilon\) 仍为一组基，故 \(A_k\simeq k[\varepsilon]/(\varepsilon^2)\)。元素 \(a+b\varepsilon\) 在 \(a\ne0\) 时有逆 \(a^{-1}-ba^{-2}\varepsilon\)，而 \(a=0\) 时不可逆。因此 \(k\varepsilon\) 是唯一极大理想，也就是 Jacobson 根。幂等条件给出 \(a^2=a\)、\(b=0\)，所以幂等元仅有零与一，其中一为本原。此时 \(C'=H+1=\varepsilon\) 非零但平方为零；规范基仍存在，代数却不半单。失效的是底域中二的可逆性，即\cref{b5:prop:hecke-rank-one-specialization}中幂等元分母 \(v+v^{-1}\) 的可逆性；在 \(v=1\) 时该分母为 \(2\)，在特征二中成为零。
\end{solution}

\begin{exercise}\label{b5:ex:A-cell-specialization}
对\subsecref{b5:subsec:A04:03}中胞腔 \(\Gamma=\{s,ts\}\) 的二维表示，沿用基 \(e_s,e_{ts}\)，计算 \(H_s,H_t\) 在 \(R=\mathbb Z[v,v^{-1}]\) 上的矩阵，并验证其中一条二次关系。再求 \(v=1\) 时两个一维胞腔的表示。
\end{exercise}
\begin{solution}
从规范基矩阵减去 \(v^{-1}I\)，得到
\[
H_s=\begin{pmatrix}v&1\\0&-v^{-1}\end{pmatrix},
\qquad H_t=\begin{pmatrix}-v^{-1}&0\\1&v\end{pmatrix}.
\]
对于 \(H_s\)，矩阵 \(H_s-vI\) 与 \(H_s+v^{-1}I\) 相乘，对角项为零，右上项为 \(0\cdot1+1\cdot0=0\)，其余也为零，故二次关系成立。胞腔 \(\{1\}\) 中，\(C'_s\cdot1=C'_s\) 属于被取商去掉的部分，所以 \(H_s\) 作用为 \(-v^{-1}\)，特化为符号表示。胞腔 \(\{sts\}\) 中 \(C'_s\) 作用为 \(v+v^{-1}\)，故 \(H_s\) 作用为 \(v\)，特化为平凡表示。
\end{solution}
